\documentclass[11pt]{article}
\pdfinfoomitdate=1
\pdftrailerid{}
\pdfsuppressptexinfo=15
\usepackage[T1]{fontenc}
\usepackage{lmodern,amsmath,amssymb,amsthm}
\pdfglyphtounicode{parenleftbig}{0028}
\pdfglyphtounicode{parenrightbig}{0029}
\pdfgentounicode=1
\usepackage[margin=1in]{geometry}
\usepackage{microtype,enumitem,needspace,booktabs,tabularx}
\usepackage[numbers,sort&compress]{natbib}
\usepackage[colorlinks=true,linkcolor=blue,citecolor=blue,urlcolor=blue]{hyperref}
\newcommand{\doi}[1]{doi: \href{https://doi.org/#1}{\nolinkurl{#1}}}
\newtheorem{theorem}{Theorem}[section]
\newtheorem{lemma}[theorem]{Lemma}
\newtheorem{proposition}[theorem]{Proposition}
\newtheorem{corollary}[theorem]{Corollary}
\newtheorem{claim}{Claim}
\theoremstyle{remark}
\newtheorem{remark}[theorem]{Remark}
\newcommand{\R}{\mathcal R}
\newcommand{\Rk}{\mathcal R_k}
\newcommand{\hc}{\operatorname{hc}}
\newcommand{\girth}{\operatorname{g}}
\numberwithin{equation}{section}
\numberwithin{table}{section}
\setlist{itemsep=3pt,topsep=5pt}
\hypersetup{
 pdftitle={Hamiltonian and Hamilton-connected thresholds for planar recolouring graphs},
 pdfauthor={},
 pdfsubject={Hamiltonian and Hamilton-connected recolouring graphs}
}

\title{Hamiltonian and Hamilton-connected thresholds for planar recolouring graphs}
\author{}
\date{}
\makeatletter
\renewcommand{\@maketitle}{%
 \newpage\null\vskip 1em
 \begin{center}
 {\LARGE \@title\par}
 \ifx\@author\@empty\else
 \vskip 1em{\large\begin{tabular}[t]{c}\@author\end{tabular}\par}
 \fi
 \ifx\@date\@empty\else\vskip .7em{\large\@date\par}\fi
 \end{center}\vskip 1em}
\makeatother

\begin{document}
\maketitle

\begin{abstract}
For a graph $G$, let $\Rk(G)$ have the proper colourings of $G$ with a
labelled $k$-colour palette as vertices, with one-vertex recolourings as
edges. We prove that $\Rk(G)$ is Hamilton-connected and vertex-pancyclic
whenever $G$ is nonempty and $d$-degenerate and $k\ge d+3$; this uniform
palette bound is sharp for both properties. We determine the common
Hamiltonian and Hamilton-connected thresholds as $7$ for planar graphs,
$5$ for triangle-free and bipartite planar graphs, and $4$ for planar
graphs of girth at least six and for forests. At four colours we also
prove Hamiltonicity for planar graphs of girth at least five with at
most thirty-one distinct $5$-cycles, total degree excess above three
at most nine, or at most seven vertices of degree at least four. We
also cover those whose blocks all have maximum degree at most three,
and two-degenerate planar graphs of girth at least five with at most
thirty-three distinct $5$-cycles or total degree excess at most thirteen.
The block condition allows unbounded degree at cutvertices.
These four-colour proofs combine structural
reductions with local restoration certificates supplied in a separate,
reproducible supplement. For unrestricted planar graphs of girth at
least five, both thresholds remain between $4$ and $5$.
\end{abstract}

\section{Introduction}

A cyclic listing of all proper colourings in which successive colourings
differ at one vertex is a colouring Gray code, or equivalently a Hamilton
cycle in the recolouring graph. Prescribing the first and last
colourings asks for Hamilton-connectedness. Requiring a cycle of every
length through each colouring asks for vertex-pancyclicity. We study
these stronger spanning and cyclic properties for sparse graphs, with
particular emphasis on sharp palette thresholds for planar classes.

All graphs are finite, simple and undirected, and may be disconnected.
A graph is nonempty when it has at least one vertex. Put $[k]=\{1,\ldots,k\}$.
The vertices of $\Rk(G)$ are the proper maps $\sigma:V(G)\to[k]$;
two are adjacent when they differ on exactly one vertex of $G$.
Colours are labelled, and a colouring need not use every colour.
A proper colouring is called \emph{frozen} when every vertex sees all
the colours other than its own on its neighbourhood; such a colouring
is an isolated vertex of the recolouring graph.
A cycle has at least three vertices. A graph is Hamilton-connected if
every two distinct vertices are joined by a spanning path, and
vertex-pancyclic if every vertex lies on a cycle of every length from
three to the order of the graph.

For a class $\mathcal C$ of nonempty graphs, define
\begin{align*}
 h(\mathcal C)&=\min\{q:\Rk(G)\text{ is Hamiltonian for every }G\in\mathcal C
                      \text{ and every }k\ge q\},\\
 \hc(\mathcal C)&=\min\{q:\Rk(G)\text{ is Hamilton-connected for every }G\in\mathcal C
                      \text{ and every }k\ge q\},
\end{align*}
when these sets are nonempty. These are uniform thresholds in the
palette size; the proofs establish every larger palette separately.
A graph is $d$-degenerate if each nonempty subgraph has a vertex of
degree at most $d$. Write $\mathcal D_d$ for this class.

Choo and MacGillivray~\cite[Theorem~3.4]{ChMa11} proved that
$\Rk(G)$ is Hamiltonian for $k\ge d+3$ when $G$ is $d$-degenerate.
They also determined the uniform forest threshold and several
individual tree and cycle cases~\cite[Sections~5--6]{ChMa11}.
Our first result strengthens their general Hamiltonian conclusion
while keeping the same palette bound.

\begin{theorem}\label{thm:main-degenerate}
Let $G$ be a nonempty $d$-degenerate graph. For every $k\ge d+3$,
$\Rk(G)$ is Hamilton-connected and vertex-pancyclic. Moreover,
$\hc(\mathcal D_d)=d+3$, and $d+3$ is also the least uniform palette
bound for vertex-pancyclicity on $\mathcal D_d$.
\end{theorem}

Write $\mathcal P$ for the planar graphs, $\mathcal T$ for the
triangle-free planar graphs, $\mathcal B$ for the bipartite planar
graphs, $\mathcal P_g$ for the planar graphs of girth at least $g$,
and $\mathcal F$ for the forests. All these classes contain only
nonempty graphs; forests have infinite girth.

\begin{theorem}\label{thm:main-planar}
The common thresholds $h=\hc$ are $7$ for $\mathcal P$,
$5$ for $\mathcal T$ and $\mathcal B$, and $4$ for
$\mathcal P_6$ and $\mathcal F$.
\end{theorem}

The large-palette cases follow from degeneracy. The critical palettes
require endpoint-preserving restoration of separated low-degree
vertices, and, in the triangle-free case, of an adjacent degree-three
pair. The lower bounds use frozen colourings and the three-colour
recolouring graph of a short path. The equality $h(\mathcal F)=4$
already follows from~\cite{ChMa11}; here it accompanies the
Hamilton-connected forest result.

At four colours the girth-five case requires different local
configurations. Write
\[
 D(G)=\sum_{v\in V(G)}\max\{d_G(v)-3,0\}
\]
for the total degree excess above three, and call a vertex heavy when
its degree is at least four. Let $\mathcal P_5^{\le31}$ consist of
the graphs in $\mathcal P_5$ with at most thirty-one distinct
$5$-cycles. Let
$\mathcal S_5$ consist of those in $\mathcal P_5$ whose every block
has maximum degree at most three. Here blocks include bridges and
isolated vertices as well as maximal two-connected subgraphs.

\begin{theorem}\label{thm:main-four}
For every $G\in\mathcal P_5^{\le31}\cup\mathcal S_5$ and every $k\ge4$,
$\Rk(G)$ is Hamiltonian. The same holds for $G\in\mathcal P_5$ with
$D(G)\le9$ or at most seven heavy vertices, and for two-degenerate
$G\in\mathcal P_5$ with at most thirty-three distinct $5$-cycles or
$D(G)\le13$. Each stated class has uniform Hamiltonian threshold four
and Hamilton-connected threshold between four and five.
\end{theorem}

In particular, Theorem~\ref{thm:main-four} includes all subcubic planar
graphs of girth at least five, while allowing vertices of arbitrarily
large total degree when that degree is distributed among blocks.
Bartier et al.~\cite{BBFHMP23} proved that $\R_4(G)$ is connected for
every $G\in\mathcal P_5$. Connectivity supplies paths between states but
does not supply a spanning cycle or prescribed endpoints, so these questions
remain strictly stronger. For the full class we obtain
\[
4\le h(\mathcal P_5)\le5,\qquad 4\le\hc(\mathcal P_5)\le5.
\]
Table~\ref{tab:thresholds} records the results and these open cases.

\begin{table}[htbp]
\centering\small
\caption{Uniform thresholds. An interval records an unresolved value.}
\label{tab:thresholds}
\begin{tabularx}{0.7\linewidth}{@{}Xcc@{}}
\toprule
Graph class & $h$ & $\hc$\\
\midrule
Planar & $7$ & $7$\\
Triangle-free planar & $5$ & $5$\\
Bipartite planar & $5$ & $5$\\
Planar, girth at least six & $4$ & $4$\\
Forests & $4$ & $4$\\
Planar, girth at least five, at most thirty-one $5$-cycles & $4$ & $[4,5]$\\
Planar, girth at least five, every block subcubic & $4$ & $[4,5]$\\
Planar, girth at least five, $D(G)\le9$ & $4$ & $[4,5]$\\
Planar, girth at least five, at most seven heavy vertices & $4$ & $[4,5]$\\
Two-degenerate planar, girth at least five, at most thirty-three $5$-cycles & $4$ & $[4,5]$\\
Two-degenerate planar, girth at least five, $D(G)\le13$ & $4$ & $[4,5]$\\
Planar, girth at least five & $[4,5]$ & $[4,5]$\\
\bottomrule
\end{tabularx}
\end{table}

\subsection{Proof strategy and organization}
Deleting vertices partitions the larger recolouring graph into fibres
above colourings of the smaller graph. At the degenerate bound these
fibres are cliques of order at least three. We choose bridge colours
between consecutive fibres to prescribe endpoints or cycle lengths.
Section~\ref{sec:degenerate} proves Theorem~\ref{thm:main-degenerate},
including the parity obstruction that makes its bound sharp.
Section~\ref{sec:restoration} develops the product and independent-set
lifts used in the planar proofs.

Section~\ref{sec:hc-planar} proves Theorem~\ref{thm:main-planar} by
combining these lifts with planar degree counts. Section~\ref{sec:four-colour-subclasses}
proves Theorem~\ref{thm:main-four}. The few-cycle argument combines the
heavy-vertex structure with restoration of low-degree configurations, including
paths whose endpoints share an outside neighbour and cycles with an attached
degree-two vertex. A degree-two vertex also supplies a reservoir for an
ordinary pentagon whose local fibre can be disconnected. Linked ordinary
pentagons and protected products for bowties and hinged pairs supply
further restorations. The new products reduce the degree-four clean-face
capacity to one. A neighbourhood face count then finishes the
minimum-degree-three case at total excess nine. For the subcubic and block-subcubic classes,
we restore two separated pentagons. Their product fibres have at most
four components. Local cycles retain protected edges at each interface;
a bounded-degree spanning tree then joins all component cycles into
one spanning cycle, for a bottom recolouring cycle of arbitrary length.

The supplement \emph{Finite certificates for Hamiltonian thresholds in
planar recolouring graphs} contains the literal paths, cycles, pools
and interface records. Its Appendices~A--J contain the finite data for these
local restoration steps. Section~\ref{sec:verification} specifies exactly
what the accompanying verifiers check. The certificate contract is deliberately
finite: the readers certify (i) legality, coverage and adjacency of each
listed local path or cycle, (ii) the stated symmetry orbits, (iii) the listed
common interface edges, and (iv) the pool and protected-edge conditions in
the finite component types. They do not certify the existence of a bottom
Hamilton cycle, an arbitrary-length bad-run lift, a product cycle for an
unlisted order, or a spanning-tree splice. Those steps are analytic
arguments in the article, and the distinction is part of the claimed
evidence rather than a hidden computational premise.

\subsection{Notation}
For a graph $G$, write $N_G(v)$ for the open neighbourhood,
$N_G[v]=N_G(v)\cup\{v\}$, and $G-S$ for deletion of a vertex set $S$.
The Cartesian product $X\square Y$ joins $(x,y)$ to $(x',y')$ when
exactly one coordinate changes along an edge of its factor.
A spanning path may be a singleton when its ends coincide in an
intermediate construction. Local certificates use zero-based palettes
$\{0,\ldots,k-1\}$, obtained by relabelling $[k]$.
Face lengths count edge-side incidences, so a bridge is counted twice.

\section{Hamilton-connectedness and vertex-pancyclicity}
\label{sec:degenerate}

We begin with the general result that supplies all large-palette statements.

\begin{lemma}
\label{lem:list-cycle}
Let $m\ge3$ and let $L_1,\ldots,L_m$ be finite sets of size at least two.  Unless $m$ is odd and all $L_i$ are the same two-element set, there are $x_i\in L_i$ with $x_i\ne x_{i+1}$ cyclically.
\end{lemma}

\begin{proof}
If some $L_i$ has at least three elements, choose successively along the path obtained by deleting that position, then choose the remaining element to avoid both end values.  If all sets are binary and two consecutive sets differ, start with a value in their symmetric difference and continue greedily.  If all are the same binary set, alternation is possible exactly when $m$ is even.
\end{proof}

\subsection{Clique fibres}
For a vertex $v$ with $H=G-v\ne\varnothing$ and $k\ge d_G(v)+3$, put
\[
 L(f)=[k]\setminus f(N_G(v)),\qquad
 F_f=\{f(c):c\in L(f)\},
\]
where $f(c)$ extends a proper colouring $f$ of $H$ by assigning colour
$c$ to $v$. The sets $F_f$ partition $V(\Rk(G))$ into cliques, and
adjacent base colourings $f,g$ satisfy
\begin{equation}\label{eq:clique-fibres}
 |L(f)|\ge3,\qquad |L(f)\cap L(g)|\ge |L(f)|-1\ge2.
\end{equation}
The cross edges are exactly $f(c)g(c)$ for $c\in L(f)\cap L(g)$.
If the base move changes a nonneighbour of $v$, then $L(f)=L(g)$.
We use these fibres throughout the single-vertex lifts.

\begin{lemma}
\label{lem:single-lift}
Let $H=G-v\ne\varnothing$ and $k\ge d_G(v)+3$. If
$f_0,f_1,\ldots,f_{s-1},f_0$ is a cycle of $\Rk(H)$ and $x\in F_{f_0}$,
then $\Rk(G)$ has a cycle through $x$ of every length in
\begin{equation}\label{eq:closed-lift-lengths}
 [2s+1,\ \textstyle\sum_i|F_{f_i}|].
\end{equation}
In particular, if $\Rk(H)$ is Hamiltonian, then so is $\Rk(G)$.
\end{lemma}

\begin{proof}
Put $S_i=L(f_i)\cap L(f_{i+1})$, with cyclic indices.
Lemma~\ref{lem:list-cycle} selects bridge colours $y_i\in S_i$
with $y_{i-1}\ne y_i$, unless $s$ is odd and all $S_i$ equal
one binary set $\{p,q\}$.

We exclude this exception. By~\eqref{eq:clique-fibres}, each
$L(f_i)$ then has size three, and
\[
 k-3=|f_i(N_G(v))|\le d_G(v)\le k-3.
\]
Thus the $k-3$ neighbours of $v$ have distinct colours, none uses
$p$ or $q$, and every base move changes a neighbour of $v$;
a nonneighbour move would leave an intersection of size three.
Order the neighbours $u_1,\ldots,u_{k-3}$. In the $(k-2)$-element
palette $[k]\setminus\{p,q\}$, append the unique unused colour
$r_i$ to the tuple $(f_i(u_1),\ldots,f_i(u_{k-3}))$.
A base move exchanges one neighbour colour with $r_i$, reversing
permutation parity. An odd number of such exchanges cannot return
to the initial tuple. Hence the bridge colours exist.

Choose an $m_0$-vertex path in $F_{f_0}$ between its bridge states
and containing $x$, with $3\le m_0\le|F_{f_0}|$.
In every other clique choose a bridge-to-bridge path on
$2\le m_i\le|F_{f_i}|$ vertices. The sums of these integer intervals
are exactly~\eqref{eq:closed-lift-lengths}. The cross edges with $v$
fixed at $y_i$ close these paths into the desired cycle. A Hamilton
base cycle with every $m_i=|F_{f_i}|$ gives a Hamilton cycle.
\end{proof}

\begin{lemma}\label{lem:bridge-colours}
Let $t\ge1$, and let $S_0,\ldots,S_{t-1}$ be finite sets of colours
such that $|S_i|\ge2$ for every $i$. Suppose that $|S_j|\ge3$ for
some $j\in\{0,\ldots,t-1\}$. For any colours $a$ and $b$, there exist
colours $x_i\in S_i$ satisfying
\[
x_0\ne a,\qquad
x_i\ne x_{i-1}\quad(1\le i\le t-1),\qquad
x_{t-1}\ne b.
\]
\end{lemma}

\begin{proof}
Set $x_{-1}=a$ and $x_t=b$. Choose $x_0,\ldots,x_{j-1}$ forwards,
avoiding the preceding colour, and $x_{t-1},\ldots,x_{j+1}$
backwards, avoiding the following colour. Each choice excludes only
one colour. Finally choose $x_j\in S_j\setminus\{x_{j-1},x_{j+1}\}$,
which is possible since $|S_j|\ge3$. The boundary assignments also
cover $j=0$, $j=t-1$, and $t=1$.
\end{proof}

\begin{lemma}\label{lem:clique-endpoints}
Let $H=G-v\ne\varnothing$ and $k\ge d_G(v)+3$.
Suppose $Y=\Rk(H)$ is nonempty and Hamilton-connected, and $v$ has
a nonneighbour in $H$. If for every $f\in V(Y)$ and $c\in L(f)$
there is a neighbour $g$ of $f$ with $c\in L(g)$, then
$\Rk(G)$ is Hamilton-connected.
\end{lemma}
\begin{proof}
Fix a nonneighbour $w$ of $v$. Colour permutations give all $k$
colours at $w$ among the states of $Y$; deleting at most one state
leaves at least $k-1\ge2$ colours. Thus every base path missing at
most one state contains a move at $w$. Its consecutive intersections
$S_i=L(f_i)\cap L(f_{i+1})$ have size at least two, and the move at
$w$ supplies one of size at least three.
Lemma~\ref{lem:bridge-colours} now gives bridge colours avoiding any
prescribed first and last colours, with distinct entrance and exit
colours in every clique. Hamilton paths inside the cliques concatenate
to a spanning path through those fibres with the prescribed endpoints.

Now prescribe distinct endpoints $\alpha(a)$ and $\beta(b)$ in
$\Rk(G)$. If $\alpha\ne\beta$, lift a Hamilton path in $Y$ between
their restrictions.

If $\alpha=\beta$, choose $c\in L(\alpha)\setminus\{a,b\}$ and
a neighbour $g$ preserving $c$. Since $|V(Y)|\ge k\ge3$, closing
a Hamilton path between the ends of $\alpha g$ gives a Hamilton
cycle through that edge. Let $h$ be its other neighbour of $\alpha$
and choose $q\in(L(\alpha)\cap L(h))\setminus\{c\}$.
Since $q$ cannot equal both $a$ and $b$, and a spanning path may be
reversed, we may interchange $a$ and $b$ so that $q\ne a$.

Partition $F_\alpha$ into paths
$P:\alpha(a)\leadsto\alpha(c)$ and
$P':\alpha(q)\leadsto\alpha(b)$.
Use a singleton for $P'$ if $q=b$, and its endpoint edge otherwise;
put all remaining vertices in $P$. Deleting $\alpha$ from the chosen
base cycle leaves a path missing only that state. Lift it from $g(c)$
to $h(q)$ through all other fibres. The cross edges at $c,q$ join
this lifted path to $P,P'$, giving the prescribed spanning path.
\end{proof}

\begin{theorem}\label{thm:degenerate-hc}
Let $G$ be a nonempty $d$-degenerate graph and $k\ge d+3$.
Then $\Rk(G)$ is Hamilton-connected.
\end{theorem}
\begin{proof}
Induct on $n=|V(G)|$, starting with $\Rk(K_1)=K_k$.
If $G=K_n$, then $\Rk(G)$ is the arrangement graph $A_{k,n}$.
Here $A_{k,n}$ is the graph of ordered $n$-tuples of distinct elements of
$[k]$, with one-coordinate changes; this is the parameter convention used
in~\cite{HLTH04}.  Since $k\ge n+2$, it is Hamilton-connected
by~\cite[Theorem~1]{HLTH04}.

Otherwise a minimum-degree vertex $v$ has degree at most $d$ and
a nonneighbour. Put $H=G-v$; induction makes $\Rk(H)$
Hamilton-connected. For each base colouring $f$ and $c\in L(f)$,
choose a vertex of degree at most $d$ in $H$. It has at least three
available colours, so can be recoloured avoiding its present colour
and $c$. This base move preserves $c$ at $v$.
Lemma~\ref{lem:clique-endpoints} applies.
\end{proof}

\begin{theorem}\label{thm:degenerate-vp}
Let $G$ be a nonempty $d$-degenerate graph and $k\ge d+3$.
Then $\Rk(G)$ is vertex-pancyclic.
\end{theorem}

\begin{proof}
Induct on $|V(G)|$, starting with $\Rk(K_1)=K_k$. For larger $G$,
choose $v$ with $d_G(v)\le d$, put $H=G-v$, use the clique fibres
above, and write $Y=\Rk(H)$, $M=|V(Y)|$, and $N=|V(\Rk(G))|$.
Colour permutations give $M\ge k\ge3$.
Fix a state $x=\alpha(a)$ of $\Rk(G)$.

The clique $F_\alpha$ already supplies a triangle through $x$.
For lengths four, five and six, recolour a vertex of degree at most
$d$ in $H$ to a colour avoiding its present colour, its neighbours,
and $a$. At most $d+2$ colours are excluded, so this gives a neighbour
$\beta$ of $\alpha$ with $a\in L(\beta)$.
Choose $p\in(L(\alpha)\cap L(\beta))\setminus\{a\}$,
$q\in L(\alpha)\setminus\{a,p\}$, and
$r\in L(\beta)\setminus\{a,p\}$.
The following cyclic sequences have the required lengths:
\[
\begin{array}{ll}
4:&\alpha(a),\alpha(p),\beta(p),\beta(a),\\
5:&\alpha(a),\alpha(q),\alpha(p),\beta(p),\beta(a),\\
6:&\alpha(a),\alpha(q),\alpha(p),\beta(p),\beta(r),\beta(a).
\end{array}
\]

For each $3\le s\le M$, induction supplies a base $s$-cycle $C_s$
through $\alpha$. Put $T_s=\sum_{f\in V(C_s)}|F_f|$.
Lemma~\ref{lem:single-lift} gives every length in $[2s+1,T_s]$.
These intervals overlap successively, since $T_s\ge3s\ge2s+3$
for $s\ge3$. Also $T_M=N$. They therefore cover every length
from seven to $N$, completing the proof.
\end{proof}

\begin{remark}\label{rem:not-edge-pancyclic}
Vertex-pancyclicity here cannot be strengthened to edge-pancyclicity
(cycles of every length through each edge). For any $k\ge4$, colour the
leaves of the $1$-degenerate star $G=K_{1,k-2}$ distinctly with
$3,\ldots,k$. The extensions with centre colours $1$ and $2$ form an
edge of $\Rk(G)$ in no triangle: such a triangle would recolour only
the centre, which has no third available colour.
\end{remark}


\begin{proposition}\label{prop:degenerate-sharp}
For every $d\ge0$, $\hc(\mathcal D_d)=d+3$. The same uniform palette
bound for vertex-pancyclicity cannot be reduced.
\end{proposition}
\begin{proof}
The upper bounds are the preceding theorems. For $d\ge1$, take
$G=K_{d+1}$ and $k=d+2$. Append the unique unused colour to each
ordered colouring of $G$. The resulting permutation of $d+2$ colours
changes by a transposition at every recolouring. Thus $\R_{d+2}(G)$
is bipartite by permutation parity, with equal parts and even order
$(d+2)!\ge6$. A spanning path must have endpoints in opposite parts,
so the graph is not Hamilton-connected. It has no triangle, so it is
not vertex-pancyclic. For $d=0$, the two-isolated-vertex graph has
$\R_2(G)=C_4$, which fails both properties.
\end{proof}

\section{Endpoint lifting and restoration}
\label{sec:restoration}

A colouring of a deleted graph determines the available colours at the
vertices to be restored.  The resulting extension graph is a fibre.  In
all the constructions below, assignments retain their actual coordinate
values: equal assignments in two different fibres are distinct vertices,
and a cross edge joins their two copies.  In particular, a move of one
vertex of the deleted graph may change several available lists at once.

\begin{lemma}
\label{lem:product}
Let $s\ge1$ and $F=\square_{j=1}^s K_{A_j}$, where every $A_j$ is
finite and $|A_j|\ge2$.  If all lists have size two, $F$ has a
Hamilton path between two prescribed vertices if and only if those
vertices lie in opposite bipartition classes.  If at least one list has
size at least three, $F$ is Hamilton-connected.
\end{lemma}

\begin{proof}
In the binary case $F$ is the hypercube $Q_s$, whose bipartition
classes have equal size.  The endpoint condition is therefore necessary.
For sufficiency, induct on $s$, starting with the single edge $Q_1$.
Choose a coordinate on which the prescribed endpoints differ and split
along it.  Their projections into $Q_{s-1}$ have the same parity.  Choose
a projected vertex of the opposite parity; use the two induction paths
to and from its copies, joined by the edge between those copies.

For the second assertion, we first show that $Y\square K_r$ is
Hamilton-connected whenever $Y$ is Hamilton-connected and
$|V(Y)|\ge3$.  Induct on $r$, starting with $r=1$.  If the two
endpoints lie in different copies of $Y$, order the copies starting and
ending at those copies.  Choose successive bridge vertices in $V(Y)$,
with the first different from the first endpoint projection, consecutive
bridges different, and the last different from the last endpoint
projection.  Such choices exist because $|V(Y)|\ge3$: choose greedily,
leaving the last bridge until the end, when it has at most two forbidden
values.  Prescribed-end Hamilton paths inside the copies then join into
a spanning path.  If both endpoints lie in one copy, take a Hamilton
path between them there and choose an edge $xy$ of it.  By induction the
other $r-1$ copies have a spanning path from a copy of $x$ to a copy of
$y$ in the same one of those copies.  Insert that path in place of
$xy$, using the two cross edges.  Start with a complete-graph factor of
order at least three and add the remaining factors to obtain the claim.
\end{proof}

\begin{lemma}
\label{lem:path-fibre}
Let $m\ge3$, $s\ge1$, and let
\[
 F_i=\square_{j=1}^s K_{A_{i,j}},\qquad |A_{i,j}|\ge2,
 \qquad 0\le i<m,
\]
be disjoint fibres arranged along a path.  Join equal assignments in
consecutive fibres.  At every transition assume that at least one list
is unchanged and that, for every coordinate,
\[
 |A_{i,j}\setminus A_{i+1,j}|\le1,
 \qquad |A_{i+1,j}\setminus A_{i,j}|\le1.
\]
If an internal fibre $F_r$, $1\le r\le m-2$, has a factor of order
at least three, then every prescribed $p\in F_0$ and $q\in F_{m-1}$
are joined by a path covering all vertices of all fibres exactly once.
\end{lemma}

\begin{proof}
Write
\[
 J_i=\square_{j=1}^s K_{A_{i,j}\cap A_{i+1,j}}
\]
for the graph of common assignments.  The intersections are nonempty,
and an unchanged coordinate supplies an edge of $J_i$.

Starting at $p$, record the reachable exit assignments at each
interface.  If a fibre has a reachable entrance $a$, an interface edge
has an admissible Hamilton-path endpoint: in a binary fibre choose its
endpoint of parity opposite to $a$, and in a nonbinary fibre choose an
endpoint different from $a$.  Lemma~\ref{lem:product} therefore shows
that the reachable set never becomes empty.

At $F_r$, its outgoing interface $J_r$ still contains an edge after any
one assignment is removed.  Indeed, if its large coordinate is
unchanged, $J_r$ contains a triangle.  Otherwise that coordinate has
intersection size at least two, and a different unchanged coordinate
supplies a $K_2\square K_2$.  Take a reachable entrance $a$ in $F_r$.
Hamilton-connectedness of $F_r$ makes every assignment in
$V(J_r)\setminus\{a\}$ reachable; this set contains an entire edge.

An incoming reachable edge permits every outgoing assignment to be
reached in the next fibre.  In a binary fibre the two endpoints of the
edge have opposite parities, so one has the parity needed for any
specified outgoing assignment.  In a nonbinary fibre they are distinct,
so one differs from that assignment.  Thus the reachable exits contain
an edge at every subsequent interface.  At the last fibre, one endpoint
of the incoming edge can be joined to the specified $q$.  Backtracking
the finite choices supplies one Hamilton path in each fibre; the actual
cross edges concatenate them into the required spanning path.
\end{proof}

\begin{lemma}
\label{lem:independent-hc}
Let $S$ be a nonempty independent set, $H=G-S\ne\varnothing$, and
$k\ge4$.  Suppose
\[
 d_G(u)\le k-2\quad(u\in S),\qquad
 |N_G(x)\cap S|\le |S|-1\quad(x\in V(H)).
\]
Assume that $\Rk(H)$ is Hamilton-connected and that some proper
$k$-colouring of $H$ leaves at least three colours available at a vertex
of $S$.  If $\alpha,\beta$ are proper $k$-colourings of $G$ with
$\alpha|_H\ne\beta|_H$, then $\Rk(G)$ has an
$\alpha$--$\beta$ Hamilton path.
\end{lemma}

\begin{proof}
Take a Hamilton path of $\Rk(H)$ between the two restriction
colourings.  For a base colouring $f$, the fibre is exactly
\[
 \square_{u\in S}K_{A_f(u)},\qquad
 A_f(u)=[k]\setminus f(N_G(u)).
\]
Every list is at least binary.  A move at an actual vertex $x$ of $H$
leaves the list of at least one non-neighbour in $S$ unchanged; every
other list loses and gains at most one colour.  These are the transition
conditions of Lemma~\ref{lem:path-fibre}, even when several lists
change together.

A colouring with a nonbinary fibre has at least $k$ distinct colour
permutations: fix a vertex of the nonempty $H$ and send its colour to
each of the $k$ labels.  Every resulting fibre is nonbinary.  The base
Hamilton path has only two endpoints, so one of these colourings is
internal.  Lemma~\ref{lem:path-fibre} now lifts the path with the
specified, actual colours of $\alpha$ and $\beta$ on $S$.
\end{proof}

For use of this lemma, first choose a vertex $w$ on which two prescribed
colourings differ, and then choose $S$ with $w\notin S$.  This ensures
different restriction colourings; the lemma does not claim to handle
arbitrary endpoints in the same fibre.

\begin{lemma}
\label{lem:fibre-lift}
Let $m\ge3$ disjoint fibres
$F_i=\square_{j=1}^sK_{A_{i,j}}$, with $s\ge1$ and
$|A_{i,j}|\ge2$, be cyclically ordered.  Join equal assignments in
consecutive fibres and define
$J_i=\square_j K_{A_{i,j}\cap A_{i+1,j}}$, with indices modulo $m$.
Suppose every $J_i$ contains an edge, $F_0$ has a factor of order at
least three, and there are distinct edges
$e_0\in E(J_0)$ and $e_{m-1}\in E(J_{m-1})$ (compared as actual
assignment pairs).  Then the layered graph has a Hamilton cycle.
\end{lemma}

\begin{proof}
Choose an edge $e_i$ in every remaining interface, and choose
$z_0\in V(e_0)\setminus V(e_{m-1})$.  For $1\le i<m$, there is a
bijection $\phi_i:V(e_{i-1})\to V(e_i)$ such that each
$z,\phi_i(z)$ is an admissible pair of Hamilton-path endpoints in
$F_i$.  In a binary fibre match opposite parities.  Otherwise use a
bijection with no fixed assignment, which always exists between two
two-element sets: interchange equal sets, send a common member to the
other target member when there is exactly one common member, and use
either bijection when the sets are disjoint.

Put $z_i=\phi_i(z_{i-1})$.  Traverse $F_i$ from $z_{i-1}$ to $z_i$
for $1\le i<m$.  Since $z_{m-1}\in V(e_{m-1})$,
$z_{m-1}\ne z_0$, so the nonbinary $F_0$ has a Hamilton path from
$z_{m-1}$ to $z_0$.  Join the fibre paths by the cross edges at the
assignments $z_i$.  They cover every fibre once and form one cycle.
\end{proof}

\begin{lemma}
\label{lem:independent-restoration}
Let $S$ be a nonempty independent set, $H=G-S\ne\varnothing$, and
$k\ge4$.  Suppose
\[
 d_G(u)\le k-2\quad(u\in S),\qquad
 |N_G(x)\cap S|\le |S|-1\quad(x\in V(H)).
\]
If $\Rk(H)$ is Hamiltonian and some proper $k$-colouring of $H$ leaves
at least three colours available at a vertex of $S$, then $\Rk(G)$ is
Hamiltonian.
\end{lemma}

\begin{proof}
Take a Hamilton cycle of the base and use the same product fibres as in
Lemma~\ref{lem:independent-hc}.  Every interface contains an edge.
Index a nonbinary fibre as $F_0$ and choose a large coordinate in it.
At either adjacent interface, that coordinate supplies a triangle if
unchanged, or a $K_2\square K_2$ together with a different unchanged
coordinate if it changes.  Thus both adjacent interfaces have at least
two edges, and distinct edges $e_0,e_{m-1}$ can be chosen.  Apply
Lemma~\ref{lem:fibre-lift}.
\end{proof}

\begin{lemma}\label{lem:single-hc}
Let $H=G-v\ne\varnothing$ and $k\ge d_G(v)+3$.
If $v$ has a nonneighbour, $\Rk(H)$ is Hamilton-connected, and
$\delta(\Rk(H))>d_G(v)$, then $\Rk(G)$ is Hamilton-connected.
\end{lemma}
\begin{proof}
For each base colouring $f$ and $c\in L(f)$, at most $d_G(v)$
moves make $c$ unavailable at $v$: only a neighbour of $v$
recoloured to $c$ can do so. The state-degree bound therefore
supplies a neighbour preserving $c$.
Apply Lemma~\ref{lem:clique-endpoints}.
\end{proof}

A Hamilton-connected graph of order at least four has minimum degree at
least three.  A vertex of degree two would, with its neighbours chosen
as the two endpoints, force a spanning path to have only three vertices;
degrees zero and one are excluded as well.  Every nonempty colourable
$H$ has at least $k$ states in $\Rk(H)$, by colour permutation.
These observations will supply the state-degree condition in the
low-degree inductions.

\section{Hamilton-connected planar classes}
\label{sec:hc-planar}

\subsection{All planar graphs}

We first establish the marked structural statement needed to retain a
coordinate on which the prescribed endpoint colourings differ.

\begin{lemma}
\label{lem:marked-five}
If $G$ is planar and $\delta(G)\ge5$, then, for any $w\in V(G)$,
there are two degree-five vertices different from $w$ whose distance is
at least three.  Vertices in different components have infinite distance.
\end{lemma}

\begin{proof}
Extend $G$, without adding vertices, to a simple plane triangulation
$T$.  A degree-five vertex in $T$ still has degree five in $G$, and a
pair at distance at least three in $T$ has that property in $G$.
Moreover $\delta(T)\ge5$.  It suffices to prove the assertion in $T$.
Euler's formula gives at least twelve degree-five vertices.

For a degree-five vertex $x$, put $W=N_T[x]$, $U=V(T)\setminus W$,
$r=|U|$, and let $B$ be the vertices at distance exactly two from $x$;
write $b=|B|$.  The five neighbours of $x$ form a cycle in their cyclic
order, so $e_T(W)\ge10$.  Contract the connected $W$ and simplify
parallel edges.  The resulting plane graph has $r+1$ vertices and
$e_T(U)+b$ edges.  As $|V(T)|\ge12$, the planar edge bound gives
$e_T(U)+b\le3r-3$.  Since $e(T)=3r+12$,
\[
 \sum_{z\in U}d_T(z)
 =e_T(U)+e(T)-e_T(W)\le6r-b-1.
\]
Let $F_x$ be the degree-five vertices at distance at least three from
$x$.  The same sum is exactly
\[
 6r-b-|F_x|
 +\sum_{z\in B}(d_T(z)-5)
 +\sum_{\substack{z\in U\setminus B\\ d_T(z)\ge6}}(d_T(z)-6).
\]
It follows that $|F_x|\ge1$, and if $|F_x|=1$, every vertex of $B$
has degree five.

Suppose every pair of degree-five vertices at distance at least three
contains $w$, and let $L$ be all degree-five vertices except $w$.
For every $x\in L$, the nonempty $F_x$ must be $\{w\}$.  Hence
$w$ is at distance at least three from $L$, and every vertex at distance
exactly two from $x$ belongs to $L$.  Take a shortest path from $w$
to $L$, ending at $x$.  The vertex two steps before $x$ is also in
$L$ and is closer to $w$, a contradiction.
\end{proof}

\begin{theorem}
\label{thm:planar-hc}
For every nonempty planar graph $G$ and every $k\ge7$, $\Rk(G)$ is
Hamilton-connected.  Consequently
$h(\mathcal P)=\hc(\mathcal P)=7$.
\end{theorem}

\begin{proof}
Prove the seven-colour assertion by induction on $|V(G)|$, allowing
disconnected graphs throughout.  If $|V(G)|\le5$, the graph is
four-degenerate and Theorem~\ref{thm:degenerate-hc} applies.
For $|V(G)|\ge6$, suppose first that some vertex $v$ has degree at
most four; it has a nonneighbour.  The smaller $H=G-v$ is Hamilton-connected by induction.
For every seven-colouring $f$ of $H$, which has at least five vertices,
\[
 d_{\R_7(H)}(f)
 =\sum_{x\in V(H)}(6-|f(N_H(x))|)
 \ge6|V(H)|-2|E(H)|\ge12>d_G(v).
\]
Lemma~\ref{lem:single-hc} restores $v$.

Otherwise $\delta(G)\ge5$.  For prescribed distinct colourings
$\alpha,\beta$, choose an actual vertex $w$ with
$\alpha(w)\ne\beta(w)$.  Lemma~\ref{lem:marked-five} supplies
$S=\{u,v\}$, avoiding $w$, with $d(u)=d(v)=5$ and
$d_G(u,v)\ge3$.  Thus $S$ is independent, no outside vertex is
adjacent to both members, and $H=G-S$ is smaller and nonempty.
Choose two nonadjacent neighbours $x,y$ of $u$; they exist since its
five neighbours cannot induce a planar $K_5$.  Properly six-colour the
planar $H-\{x,y\}$, and give both $x,y$ colour seven.  This is a
proper colouring of $H$ in which the neighbours of $u$ use at most
four colours, leaving at least three available.  The restriction
colourings differ at $w$, so Lemma~\ref{lem:independent-hc} gives
an $\alpha$--$\beta$ Hamilton path.

For $k\ge8$, planar five-degeneracy and
Theorem~\ref{thm:degenerate-hc} suffice.  For sharpness, use the following
explicit model of the icosahedron.  Put $\varphi=(1+\sqrt5)/2$ and take the
twelve points
\[
 \{(0,\pm1,\pm\varphi),(\pm1,\pm\varphi,0),
          (\pm\varphi,0,\pm1)\},
\]
joining two points at distance $2$.  The six antipodal pairs are independent,
and a direct check of the five neighbours of any point shows that they contain
exactly one point from each of the other five pairs.  Colour each antipodal
pair with its own colour.  Every vertex then sees all other five colours, so
the colouring is frozen: no one-vertex recolouring is possible.  Relabelling
the six colours gives distinct frozen states, so the recolouring graph is
disconnected.  This is the standard icosahedral obstruction (see also
\cite{BP16}), and proves that both thresholds are seven.
\end{proof}

Planar five-degeneracy and Theorem~\ref{thm:degenerate-vp} also give
vertex-pancyclicity for $k\ge8$.

\subsection{Triangle-free and bipartite planar graphs}

The critical five-colour step needs an adjacent-pair fibre, which is
not a product of complete graphs.  We give its finite input separately
from the arbitrary-size lifting argument.

For $A,B\subseteq\{0,1,2,3,4\}$, put
\[
 F(A,B)=\{(a,b)\in A\times B:a\ne b\},
\]
with adjacency given by changing one coordinate.  When $|A|,|B|$ are
three or four, exchange the coordinates if necessary so that
$|A|\le|B|$.  There are exactly seven types
$(|A|,|B|,|A\cap B|)$.

\begin{lemma}
\label{lem:pair-five-endpoints}
For each of the seven types
\[
331,332,333,342,343,443,444,
\]
every edge of $F(A,B)$ lies on a Hamilton cycle.  The type $333$ is
itself a $6$-cycle.  The types $342,343,443,444$ are
Hamilton-connected.  If $J$ is the common state set of one of the
interfaces allowed in Lemma~\ref{lem:pair-five-reset}, then every
$p\in F(A,B)$ has a neighbour $q\in J$, and hence a $p$--$q$
Hamilton path in $F(A,B)$.
\end{lemma}

\begin{proof}
Supplementary Appendix~A gives one short Hamilton cycle for each of
the seven representatives and checks its orbit under the independent
row and column permutations that agree on $A\cap B$ (with the stated
side exchange for type $332$).  The resulting edge orbits are the full
edge sets.  The 53 endpoint words in supplementary Appendix~A prove the four
Hamilton-connected types. Their symmetry coverage and an exact subset
dynamic program verify the complete endpoint classification.

For the interface exit, exchange coordinates if needed so that
$B$ is unchanged, and put $A_0=A\cap A'$. All allowed transitions,
including the unchanged case, give $|A_0|\ge2$ and $|B|\ge3$.
For $p=(a,b)$ with $a\in A_0$, choose $(a,d)$ with
$d\in B\setminus\{a,b\}$; otherwise choose $(c,b)$ with
$c\in A_0\setminus\{b\}$. Either choice is a neighbour in $J$.
Deleting this edge from a Hamilton cycle gives the required
$p$--$q$ Hamilton path.
\end{proof}


\begin{lemma}
\label{lem:pair-five-reset}
Arrange full fibres $F(A_i,B_i)$ along a base path.  Each list has
size three or four in a fixed five-colour palette.  At each step either
both lists stay unchanged, or exactly one changes. A changed list
is either replaced by another three-element list differing in one
colour, or changes between a three-element list and a containing
four-element list.
Join the copies of equal assignments at consecutive layers.  If an
internal layer has $|A_r|=|B_r|=4$, there is a spanning path between
any prescribed first-layer and last-layer states.
\end{lemma}

\begin{proof}
First consider one allowed interface between $F(A,B)$ and $F(A',B')$,
and write $J$ for their common assignment set.  The preceding lemma
supplies, in both directions, a Hamilton path from every entrance state
to some state of $J$.

Use this property to propagate nonempty sets of reachable exits from
the prescribed first state.  The internal $44$ layer is
Hamilton-connected, and its outgoing interface has at least nine states.
Any reachable entrance $p$ can therefore reach all interface states
except possibly $p$, yielding at least two distinct reachable exits.
The next layer still has one four-element list and is one of
$342,343,443,444$, hence is Hamilton-connected.  Two distinct reachable
entrances let every state of that layer be an exit.  If it is the last
layer, this includes the prescribed endpoint.  Otherwise its entire
outgoing interface is reachable.

Thereafter this property propagates: for a specified exit $q$ in the
current layer, apply the one-interface property in reverse to find a
Hamilton-path entrance in its fully reachable incoming interface.
The last prescribed endpoint is therefore reachable.  Backtracking
gives one Hamilton path in every full fibre, joined by actual cross
edges, with no repetition or omission.
\end{proof}

\begin{theorem}
\label{thm:trianglefree-hc}
If $G$ is a nonempty triangle-free planar graph, then $\Rk(G)$ is
Hamilton-connected for every $k\ge5$.  For the class of nonempty
triangle-free planar graphs,
$h=\hc=5$.
\end{theorem}

\begin{proof}
For $k\ge6$, use three-degeneracy and
Theorem~\ref{thm:degenerate-hc}.  For $k=5$, induct on the order of
$G$, permitting disconnected graphs.  Orders at most three are handled
by two-degeneracy and the same theorem.  At larger orders, suppose
first that some vertex $v$ has degree at most two; it has a nonneighbour.  The induction hypothesis
gives a Hamilton-connected base of at least five states, thus of minimum
degree at least three; Lemma~\ref{lem:single-hc} applies.

Assume now $\delta(G)\ge3$.  Let $L$ be its degree-three vertices
and $s=|L|$.  The triangle-free planar edge bound gives
\[
 s\ge8+\sum_{d_G(z)\ge4}(d_G(z)-4).
\]
Given distinct $\alpha,\beta$, choose $w$ on which they differ.
There are two exhaustive cases.

If $L\setminus\{w\}$ contains adjacent $u,v$, denote their other
neighbours by $u_1,u_2$ and $v_1,v_2$.  The four actual outside
vertices are distinct: a common neighbour of $u$ and $v$ would form
a triangle.  In $H=G-\{u,v\}$ the restriction endpoints still differ
at $w$, so induction gives a base Hamilton path between them.  Its
fibre over $f$ is
\[
 F(A_f,B_f),\qquad
 A_f=[5]\setminus\{f(u_1),f(u_2)\},\quad
 B_f=[5]\setminus\{f(v_1),f(v_2)\}.
\]
The lists have size three or four.  One actual base move affects at
most one of them and gives precisely an interface allowed in
Lemma~\ref{lem:pair-five-reset}.  Edges among the four outside
vertices restrict the base colourings, but do not alter this statement.

There is a base colouring with both lists four-element: use
Gr\"otzsch's theorem~\cite{Grotzsch59} to three-colour
$H-\{u_1,u_2,v_1,v_2\}$, then give $u_1,u_2$ colour four and
$v_1,v_2$ colour five.  Each pair is independent by triangle-freeness,
and edges across the pairs have distinct colours.  At least five global
colour permutations give distinct double-four-list base states.  Since
the Hamilton path has only two endpoints, one is internal.  Apply
Lemma~\ref{lem:pair-five-reset} with the specified local endpoint
states to complete the restoration.

Otherwise $S=L\setminus\{w\}$ is independent and has at least
$s-1\ge7$ vertices; put $H=G-S$.  For any outside vertex $z$ of degree at least
four, the displayed Euler budget implies $d_G(z)\le s-4<|S|$.
The only possible degree-three outside vertex is $w$, whose degree is
also less than $|S|$.  Thus no outside vertex sees all of $S$.
Choose $u\in S$ and two of its neighbours $x,y$.  They are
nonadjacent.  Three-colour $H-\{x,y\}$ by Gr\"otzsch's theorem
and give both $x,y$ colour five.  At most two colours then occur on the
three neighbours of $u$, so $u$ has at least three available colours.
All conditions of Lemma~\ref{lem:independent-hc} hold, and $w$ is
retained; it gives the desired endpoint path.

For sharpness, on the cube $Q_3$ view the colours as
$\mathbb Z_2^2$ and set
\[
 f(x,y,z)=x(1,0)+y(0,1)+z(1,1).
\]
The three neighbours of every vertex have all three other colours,
so the colouring is frozen and the complete four-colour recolouring
graph is disconnected.  The upper and lower bounds give both thresholds.
\end{proof}

\begin{corollary}
The class of nonempty bipartite planar graphs has $h=\hc=5$.
For every graph in this class, $\Rk(G)$ is vertex-pancyclic for
$k\ge6$.
\end{corollary}

\begin{proof}
Bipartite planar graphs are triangle-free, and the cube belongs to this
class.  The vertex-pancyclic assertion follows from three-degeneracy and
Theorem~\ref{thm:degenerate-vp}.
\end{proof}

The same vertex-pancyclic bound $k\ge6$ holds for all triangle-free
planar graphs.  This conclusion uses the degenerate theorem separately
from Hamilton-connectedness.

\subsection{Girth at least six and forests}

\begin{lemma}
\label{lem:marked-two}
If $G$ is planar, $\girth(G)\ge6$, and $\delta(G)\ge2$, then for
any $w\in V(G)$ there are two degree-two vertices different from
$w$ whose distance is at least three.
\end{lemma}

\begin{proof}
Every component has a cycle, and the girth-six Euler bound gives at
least six degree-two vertices per component.  If there are at least two
components, choose such a vertex different from $w$ in each of two
components.  Suppose $G$ is connected and let $n_2$ be its number of
degree-two vertices.  From $2|E(G)|\le3|V(G)|-6$,
\[
 n_2\ge6+\sum_{d_G(z)\ge3}(d_G(z)-3).
\]
Let $L$ consist of all degree-two vertices except $w$; then
$|L|\ge n_2-1\ge5$.  Suppose all pairs in $L$ have distance at
most two.

If $u,v\in L$ are adjacent, let $a,b$ be their respective other
neighbours.  The absence of cycles of length at most five implies that
the intersection of the two closed distance-two balls is exactly
$\{u,v,a,b\}$, and that $d_G(a,b)=3$.  Thus at most three of
these vertices can lie in $L$, a contradiction.

Otherwise $L$ is independent.  Its two-element neighbourhoods are
pairwise intersecting; they are distinct because there is no four-cycle.
At least four distinct pairwise intersecting two-element sets have a
common element: if $\{a,b\},\{a,c\}$ occur and a third set omits
$a$, it must be $\{b,c\}$, after which no fourth distinct set is
possible.  Hence some vertex $z$ is adjacent to all of $L$.
Then $d_G(z)\ge n_2-1\ge5$, whereas the Euler inequality gives
\[
 n_2\ge6+d_G(z)-3\ge n_2+2,
\]
a contradiction.
\end{proof}

\begin{theorem}
\label{thm:girth6-hc}
For every finite integer $g\ge6$, the class of nonempty planar graphs
of girth at least $g$ has $h=\hc=4$.  Every graph $G$ in this
class has Hamilton-connected $\Rk(G)$ for $k\ge4$, and
vertex-pancyclic $\Rk(G)$ for $k\ge5$.
\end{theorem}

\begin{proof}
For $k\ge5$, use two-degeneracy and
Theorems~\ref{thm:degenerate-hc} and~\ref{thm:degenerate-vp}.
At $k=4$, induct on the order, starting with orders at most two,
which are forests and covered by the one-degenerate theorem.
If $|V(G)|\ge3$ and $d_G(v)\le1$, then $v$ has a non-neighbour.
The smaller nonempty base is Hamilton-connected by induction and has
at least four states, hence minimum degree at least three.  Restore
$v$ using Lemma~\ref{lem:single-hc}.

Otherwise apply Lemma~\ref{lem:marked-two} to a vertex $w$ at which
the prescribed endpoint colourings differ.  It gives separated
degree-two vertices $u,v$, both different from $w$.  Set
$S=\{u,v\}$ and $H=G-S$.  The set $S$ is independent, no outside
vertex sees both its members, and $H$ is nonempty.  The neighbours
$x,y$ of $u$ are nonadjacent.  The two-degenerate $H-\{x,y\}$
can be three-coloured; giving $x,y$ colour four leaves three colours
available at $u$.  Lemma~\ref{lem:independent-hc} now restores a
Hamilton path between the actual prescribed states.

For the lower bound, take a cycle of length at least $g$ divisible by
three and repeat the colours $0,1,2$ around it.  This is a frozen
proper three-colouring, so its complete three-colour recolouring graph
is disconnected.  Both class thresholds are therefore four.
\end{proof}

\begin{theorem}
\label{thm:forest-hc}
The class of nonempty forests has $h=\hc=4$.  Every such forest
$F$ has Hamilton-connected and vertex-pancyclic $\Rk(F)$ for
$k\ge4$.
\end{theorem}

\begin{proof}
The upper bounds follow from one-degeneracy and
Theorems~\ref{thm:degenerate-hc} and~\ref{thm:degenerate-vp}.
For the Hamiltonian lower bound, $\R_3(P_3)$ has twelve
states.  Its six states using three distinct colours have degree two,
so every edge incident with them is forced in a spanning cycle.
For each of the three middle-vertex colours $b$, these forced edges
form a $4$-cycle
\[
 (a,b,a),(c,b,a),(c,b,c),(a,b,c),
 \qquad \{a,b,c\}=\{0,1,2\}.
\]
The three disjoint forced cycles prevent a Hamilton cycle.  For the
Hamilton-connected lower bound, $\R_3(K_2)=C_6$
is already not Hamilton-connected.
\end{proof}

\section{Hamiltonian four-colour subclasses at girth five}
\label{sec:four-colour-subclasses}

The reductions below splice Hamilton cycles and establish
Hamiltonicity at four colours. They do not prescribe arbitrary
endpoints, so the Hamilton-connected thresholds of these subclasses
remain in $[4,5]$. Supplementary Appendices~B--J supply their local
certificates.

Call a vertex \emph{heavy} if its degree is at least four, and a pentagon
\emph{ordinary} if all five of its vertices have degree three. These
degrees are measured in the graph before deletion. Vertices of degree
three are called \emph{light}.

\subsection{Graphs with few five-cycles}
\begin{lemma}
\label{lem:subdivided-star-four-subclass}
Let $S=\{u_1,\ldots,u_s\}$ be independent, $s\ge3$, and suppose
$N(u_i)=\{w,x_i\}$, where $N(w)=S$, $w\notin S$, and the $x_i$ are distinct and belong to
$J=G-(S\cup\{w\})$.  If $\R_4(J)$ is Hamiltonian and some proper colouring
of $J$ gives the same colour to two distinct $x_i$, then $\R_4(G)$ is
Hamiltonian.
\end{lemma}

\begin{proof}
Take a Hamilton cycle $C$ of $\R_4(J)$.  Write three distinct chosen edges
as $e_r=f_rg_r$ $(r=1,2,3)$.  Since one edge of $C$ changes only one vertex
of $J$ and the $x_i$ are distinct, choose $j(r)$ so that $x_{j(r)}$ is
unchanged on $e_r$, and put
$a_r=f_r(x_{j(r)})=g_r(x_{j(r)})$.

There is a tree $T$ on the four colours whose edge $t_r$ contains $a_r$.
Indeed, if the three labels are distinct, use the star centred at the fourth
colour; if exactly two labels are equal, say $a,a,b$, use the edges
$ab,ac,bd$, where $c,d$ are the two remaining colours; and if all labels
are equal, use the star centred at that common label.  Make four copies of
$C$, indexed by the colour of $w$.  For each $r$, remove the copy of $e_r$
from the two copies indexed by the ends of $t_r$, and add the two cross
edges joining corresponding $J$-colourings while changing the colour of
$w$.  These are legal edges of $\R_4(G-S)$ because $w$ is isolated in
$G-S$.  The three two-edge switches join the four cycles according to the
tree $T$, so the resulting $2$-regular connected graph is a Hamilton cycle
$B$ of $\R_4(G-S)$.

The only edges of $B$ that change the colour of $w$ are these cross edges.
At the one indexed by $r$, the unchanged vertex $x_{j(r)}$ has colour $a_r$,
and $a_r$ is one of the two endpoint colours of $t_r$.  Thus every such
transition has an unchanged $x_j$ whose colour is one of the two colours of
$w$; this is the required control property.

Over a base colouring $f$ the $u_i$-lists are
\[
 A_i(f)=[4]\setminus\{f(w),f(x_i)\}.
\]
They have size at least two.  At a move in $J$, one list changes and another
coordinate is unchanged.  At a move of $w$, all intersections are
$[4]\setminus\{a,b,f(x_i)\}$, and the control property leaves a coordinate
with a two-element intersection.  Thus every interface has an edge, and the
cyclic product-fibre lift applies once we exhibit its nonbinary fibre and the
two distinct incident interface edges.  Let $h$ be the assumed colouring of
$J$ with $h(x_p)=h(x_q)=c$, and take the base state with restriction $h$ and
$w$ coloured $c$.  The two lists at $u_p,u_q$ then both have size three.  If
an incident edge of $B$ changes a vertex of $J$, at most one of these two
lists changes, so the other unchanged list supplies a $K_3$ in the common
interface.  If it changes $w$ from $c$ to $b$, both common lists contain
$[4]\setminus\{c,b\}$ and hence supply the two distinct edges of
$K_2\square K_2$.  Consequently the two incident interfaces contain at
least two distinct edges, from which distinct edges can be chosen for the
cyclic lift.  Its fibre over the displayed base state has a factor $K_3$;
the lift therefore gives a Hamilton cycle through every extension of every
base state, proving the lemma.
\end{proof}

\begin{lemma}
\label{lem:five-face-identity-subclass}
Let $C$ be a connected plane graph of girth at least five and minimum degree
at least two.  Let $p$ be the number of $5$-faces,
$Q=\sum_{|F|\ge7}(|F|-6)$, $s=|\{v:d(v)=2\}|$, and
$D=\sum_v(d(v)-3)$ over vertices of degree at least three.  Then
\[
 p=12+Q+2D-2s.
\]
Every $5$-face is a simple $5$-cycle, and distinct faces give distinct
cycles unless $C=C_5$.
\end{lemma}

\begin{proof}
A facial walk has no immediate reversal and contains a cycle, hence has
length at least five.  Counting face lengths gives
$2e=6f-p+Q$, while the degree sum is $2e=3n-s+D$; with $n-e+f=2$,
subtraction yields the identity.  A repeated vertex on a facial walk of
length five would give a shorter cycle.  If one $5$-cycle bounds two faces,
all its vertices have degree two and connectedness gives $C=C_5$.
\end{proof}

\begin{corollary}
\label{cor:eleven-pentagons-two-degenerate}
A planar graph of girth at least five with at most eleven distinct $5$-cycles
is two-degenerate.
\end{corollary}

\begin{proof}
Otherwise a connected subgraph has minimum degree at least three.  In its
plane embedding the identity has $s=0$, and therefore at least twelve
$5$-faces, which are distinct $5$-cycles.  This contradicts the hypothesis.
\end{proof}

\subsection{Local degree-two reductions}
\begin{lemma}
\label{lem:adjacent-two-four}
Let $u,v$ be adjacent degree-two vertices, with
$N(u)=\{v,x\}$ and $N(v)=\{u,y\}$, where $x\ne y$.  If
$\R_4(G-\{u,v\})$ is Hamiltonian, then $\R_4(G)$ is Hamiltonian.
\end{lemma}
\begin{proof}
For a base colouring $f$, the extension fibre is $F(A,B)$, where
$A=[4]\setminus\{f(x)\}$ and $B=[4]\setminus\{f(y)\}$.
For $A=B=\{a,b,c\}$ use the cycle $(ab,ac,bc,ba,ca,cb)$;
for $A=\{a,b,c\}$, $B=\{a,b,d\}$ use
$(ab,ad,bd,ba,ca,cd,cb)$.

In fact, any two fibre Hamilton cycles have at least two common
edges when at most one list changes. For an unchanged fibre, equal
lists give $C_6$, while unequal lists give four distinct forced
edges at the two degree-two states using the common colours.
By exchanging coordinates suppose only $A$ changes. If one value
equals $B=\{a,b,c\}$, the common fibre is the path
$(ab,ac,bc,ba)$; its two end edges are forced in both cycles.
Otherwise let $a,b$ be the old and new excluded colours, $c$ the
fixed excluded colour and $d$ the fourth colour. The common fibre
has five vertices $ca,cb,cd,da,db$ and six edges. The old outside
states $ba,bd$ occur consecutively because $bd$ has degree two
with neighbours $ba,cd$. Thus the old cycle uses exactly four
common-fibre edges; the new cycle does likewise with $ab,ad,cd$.
The cycles share at least $4+4-6=2$ edges.

Since $x\ne y$, a base step changes at most one list. Cut an edge
of a base Hamilton cycle and process its fibres in path order.
At each splice the current fibre has lost at most one edge and the
next is fresh, so a common edge remains. Cutting its two copies
and adding the two actual base-step edges merges the circles.
The fibres partition all states, giving a spanning cycle.
\end{proof}

\begin{lemma}\label{lem:separated-two-g5}
Let $G$ be planar of girth at least five and let $u,v$ be nonadjacent
vertices of degree two with $N(u)\cap N(v)=\varnothing$.  If
$H=G-\{u,v\}$ has Hamiltonian $\R_4(H)$, then $\R_4(G)$ is Hamiltonian.
\end{lemma}
\begin{proof}
Set $S=\{u,v\}$.  Every vertex of $H$ has at most one neighbour in $S$.
Let $x,y$ be the neighbours of $u$; they are distinct and nonadjacent.
The planar triangle-free graph $H-\{x,y\}$ admits a proper three-colouring
by Gr\"otzsch's theorem~\cite{Grotzsch59} (the empty graph has its empty
colouring).  Use colours $1,2,3$ there and give both $x,y$ colour $4$.
This is a proper colouring of $H$ leaving three colours available at $u$.
The hypotheses of Lemma~\ref{lem:independent-restoration} hold with
$k=4$ and $|S|=2$, so that lemma supplies the required Hamilton cycle.
\end{proof}

\subsection{\texorpdfstring{Restoring a $(2,3,2)$ path}{Restoring a (2,3,2) path}}
The following reduction uses only a Hamilton cycle of the deleted graph.
It processes an entire constant-boundary run together with its two
neighbouring nonconstant fibres, retaining the edges needed for the next
splice.  The three external neighbours remain distinct actual vertices.
\begin{proposition}\label{prop:badrun}
If $F=\R_3(P_3)$, then $F\mathbin{\square}P_\ell$ is Hamiltonian
for every $\ell\ge2$.
\end{proposition}
\begin{proof}
Partition $F$ into its three centre-colour squares $Q_b$.
Order each square as $(q_0,q_1,q_2,q_3)$, where $q_0,q_2$ are
its monochromatic-endpoint vertices. On $Q_b\square P_2$ use
\[
 (0,q_0),(1,q_0),(1,q_1),(0,q_1),
 (0,q_2),(1,q_2),(1,q_3),(0,q_3).
\]
This cube cycle contains all four edges between layers 0 and 1.
For larger $\ell$, append each subsequent square by an ordinary
edge splice using an internal edge of the newest layer.  The induction
invariant is that the strip is one cycle, contains every vertex in the
first $\ell$ layers, retains all four edges between layers 0 and 1,
and has at least two unused internal edges in its newest layer.  The
initial strip has two such edges; after one splice the new layer has
three, so a further splice can avoid any prescribed reserved edge.
Thus the invariant holds for every $\ell\ge2$, and the four original
interlayer edges are never cut.

The three centre-colour squares of $F$ are connected by three edges
between monochromatic-endpoint vertices, one for each pair of
centre colours. Choose a tree of two of these edges. Each chosen
edge $xy$ produces a square with vertices $(0,x),(1,x),(1,y),(0,y)$.
Delete its two interlayer edges and add its two copies of $xy$,
joining the corresponding strip cycles. The two selected interlayer
edges at the middle strip are distinct. Two such splices give a
single spanning cycle, for arbitrary $\ell\ge2$.
\end{proof}

Let $x,y,z$ be a path of degrees $2,3,2$.  Its external neighbours, in
that order, are $p,q,r$.  Assume they are distinct and outside the path.
For a bottom colouring with boundary $t=(f(p),f(q),f(r))$, write
\[
 B_t=\{(u,v,w):u\ne v,\ v\ne w,\ u\ne t_0,\ v\ne t_1,
                  \ w\ne t_2\},\qquad u,v,w\in\{0,1,2,3\}.
\]
A boundary is \emph{bad} if all three entries agree, and \emph{good}
otherwise. Up to colour permutations and path reversal the good types
are $001,010,012$. Their orders are $6+4+4=14$, $9+4+4=17$, and
$6+6+4=16$, respectively. The bad type $000$ has order $12$.

\paragraph{Good fibres and safe edges.}
The cycles $D_E,D_M,D_A$ for the normalized boundaries $001,010,012$
are specified by literal cycles and capacity checks in supplementary
Appendix~B; the coordinates are $uvw$.
Allow all their colour images and path reversals. Legality, cyclic adjacency, distinctness, and the counts
above verify that each is a Hamilton cycle of its fibre.

For a bad boundary $aaa$, partition $B_{aaa}$ into squares $Q_c$ indexed
by $c\ne a$, with fixed middle coordinate $c$. A good neighbour of
$aaa$ is obtained by changing position $j$ from $a$ to $b\ne a$.
If $j$ is an end position, then for each $c\notin\{a,b\}$ there is
exactly one edge common to the good fibre and $Q_c$. Let $P_t$ consist
of these two edges. For middle-change or all-distinct boundaries set
$P_t=\varnothing$. A good-cycle edge outside $P_t$ is called \emph{safe}.
Directly from the supplementary path certificate, every allowed cycle has the following properties:
\begin{enumerate}
\item For an end change, it contains the unique edge in each
      $Q_c$, $c\notin\{a,b\}$, and two opposite edges of $Q_b$.
\item For a middle change it has no vertex in $Q_b$ and contains
      three edges of each of the other two squares.
\end{enumerate}
For $t=001$, the two special edges are $123$--$323$ and $132$--$232$;
the two edges in $Q_1$ are $313$--$213$ and $212$--$312$.
For $t=010$, the retained three-edge paths are
$121,321,323,123$ and $132,232,231,131$.
These descriptions also prove the assertions for every symmetry image.

\begin{lemma}\label{lem:pathsafe}
For adjacent good boundaries $t,s$, every allowed cycle at $t$ and
every allowed cycle at $s$ have at least two common edges outside
$P_t\cup P_s$.
\end{lemma}
\begin{proof}
Partition the common fibre graph into its fixed-middle layers $H_v$
and the edge set $H_*$ changing the middle coordinate. In any block,
if $m$ edges remain outside $P_t\cup P_s$ and the two cycles retain
$x,y$ of them, they share at least $\max(0,x+y-m)$ safe edges.
These bounds add over disjoint blocks. The supplementary path-kernel certificate
classifies all allowed transitions and gives sufficient block counts
for seven families, each yielding at least two.

Let $a,b,c,d$ be distinct and use the selected cycles defined there.
For $aab\to bab$, the common middle-$b$ prism has no special edges.
The $E$ cycle retains one edge in each row triangle and two vertical
edges. The cycle $M_0$ contains both these triangle edges, and $M_1$
both these vertical edges, giving two safe common edges in either case.

For $aba\to cba$, $H_a$ has nine edges. Both cycles omit
$bab$--$bac$ and $bab$--$bad$: the $M$ layer isolates $bab$, while
the $A$ path starts $bab,dab$. They retain three and five edges in
this seven-edge universe, giving one common edge. In $H_c$ they
retain three and two of four edges, giving another. Both special
sets are empty. Colour permutations and reversal transport all
arguments, proving the assertion for every allowed labelled pair.
\end{proof}

\paragraph{Bad blocks.}
\begin{lemma}\label{lem:badblock}
Consider consecutive fibres with boundaries
\[
             t,\underbrace{aaa,\ldots,aaa}_{\ell\text{ times}},s,
             \qquad \ell\ge1,
\]
where $t,s$ are good and each differs from $aaa$ in one coordinate.
When $\ell=1$ suppose $t\ne s$. Choose any allowed good cycles at
the two ends, and reserve at most one safe edge in each. The layered
graph on these fibres has a Hamilton cycle retaining both reserved
edges. Only consecutive-layer matching edges are used.
\end{lemma}
\begin{proof}
Call a common edge with $Q_c$ \emph{available} if it belongs to the
chosen good cycle and is not its reserved edge. An end-change good
cycle has available edges in all three squares: the two singleton
edges are special, hence cannot have been reserved, and at least one
of its two $Q_b$ edges remains. A middle-change good cycle has at least
two available edges in each of its two supported squares.

First suppose $\ell=1$. Let $S_L,S_R$ be the sets of square
colours supported by available edges of the two chosen good cycles.
Then $|S_L|,|S_R|\ge2$ and $S_L\cup S_R=[4]\setminus\{a\}$:
the only exception could be two middle changes to the same colour,
which would give $t=s$.

Choose a common square with different available edges on the two
sides. If either change is at the middle, any $c\in S_L\cap S_R$
works, since that side offers at least two edges. If both changes
are at ends, choose $c\notin\{a,b_L,b_R\}$, where $b_L,b_R$ are
their changed colours. Its two singleton edges are different:
different ends give different coordinate directions, while changes
at the same end have $b_L\ne b_R$ and give opposite square edges.

Start with the two good cycles and three bad squares. Attach the
chosen square to both good cycles using these different edges, and
attach each other square to a supporting good cycle. These four
incidences form a tree. Cuts at a good cycle lie in different
middle-colour squares, and the two bridge-square cuts are distinct.
Along each tree edge delete the two selected internal edges and add
the two matching edges. Each splice joins disjoint cycles into one;
all reserved edges survive.

Now let $\ell\ge2$. Each good end supports at least two squares, so
choose an available left edge in $Q_c$ and an available right edge in
$Q_d$ with $c\ne d$. Build the three square-strip cycles as in
Proposition~\ref{prop:badrun}. The cube formula there permits either
perfect matching of a square to be used in the first layer, by rotating
its cyclic order; the last layer of a two-layer strip uses the other
matching. Consequently its first layer can retain any specified square
edge. For a two-layer strip the same is true at its last layer; our two
requirements are on different strips and are independent.

If $\ell\ge3$, a last-layer requirement is met by appending each new
square with a splice edge different from the required edge. There are
initially two choices and subsequently three in the current last
layer, so this is always possible. The new last square retains all
edges except the splice edge. First-layer horizontal edges and all
four original vertical edges survive these extensions.

Join the three strips using the two vertical-edge switches in
Proposition~\ref{prop:badrun}; these do not affect either required
horizontal edge. Finally join the left and right good cycles to their
selected strip edges by two ordinary edge splices. The result is one
Hamilton cycle. The reserved edges were never cut. This proves the
claim for every length, rather than just a bounded set of runs.
\end{proof}

\paragraph{Restoration over a base cycle.}
\begin{lemma}\label{lem:path232}
Suppose the three external neighbours of a degree-$(2,3,2)$ path are
distinct and outside it. If the reconfiguration graph of the deleted
graph is Hamiltonian at four colours, then so is $\R_4(G)$.
\end{lemma}
\begin{proof}
Fix an arbitrary bottom Hamilton cycle. There must be a good vertex
on it. Otherwise every boundary is constant; consecutive constant
triples can only be equal, as at most one coordinate changes. The
external colour would therefore be fixed throughout the cycle. A
global colour permutation of a proper colouring supplies another
proper colouring with a different external colour, a contradiction
to Hamilton coverage. Thus even the startup uses only Hamiltonicity.

Every good bottom vertex has at most one bad neighbour in the entire
bottom reconfiguration graph: a two-colour boundary has exactly one
possible coordinate change making it constant, and an all-distinct
boundary has none. The full neighbouring colouring, if proper, is
uniquely determined by that change. Hence no good vertex on the
Hamilton cycle can lie between two bad vertices. In particular there
is an edge joining two good vertices.

Cut such an edge. In the resulting bottom Hamilton path, group each
maximal bad run together with its two neighbouring good vertices;
make each remaining good vertex a singleton block. These blocks are
disjoint and partition the path. A bad run has constant boundary,
and its two neighbours are distinct. If the run has length one their
boundaries are different: otherwise both neighbours are obtained from
the intervening colouring by the same one-vertex recolouring and
would be the same bottom vertex.

For each labelled good boundary fix one allowed cycle once and for
all. Consecutive blocks meet at a good--good boundary transition.
If the boundary changes, Lemma~\ref{lem:pathsafe} provides two common
safe edges. If it stays unchanged, the fixed cycles coincide and
have at least $14-2$ safe edges. Select one connecting edge for each
interblock interface in path order. If a singleton block already has
an incoming selected edge, choose a different outgoing one, which is
possible from the two candidates. In a nonsingleton block the two
selected edges are in its distinct end fibres, so cannot conflict.

Use Lemma~\ref{lem:badblock} to construct each nonsingleton block
cycle preserving its selected end edges; singleton blocks already
have their fixed cycles. Splice the resulting disjoint block cycles
along the selected interfaces. Each step joins two cycles into one,
and every later selected edge is still present. The one-block case
needs no final splice. All transitions are legal: a local step changes
one path vertex, and a matching cross-edge performs its actual bottom
recolouring while holding the local colouring fixed. Finally the fibres
are disjoint and exhaust all proper colourings. Thus the resulting
single cycle is Hamiltonian, including its closing edge.
\end{proof}


\subsection{Paths with a common outside neighbour}

For $r\in\{4,5\}$ and $t=(a,b_1,\ldots,b_{r-2})\in Q^{r-1}$, where
$Q=\{0,1,2,3\}$, let $F_t^{(r)}$ consist of the proper colourings of
the path $v_0\cdots v_{r-1}$ that avoid the respective colours
$(a,b_1,\ldots,b_{r-2},a)$. Two words are adjacent when they differ in
one coordinate. The palette $Q$ is a relabelling of $[4]$; the first
boundary coordinate is the colour of one outside neighbour shared by
the endpoints.

\begin{lemma}\label{lem:shared-hub-path-data}
There are Hamilton cycles $C_t^{(r)}$ in all $F_t^{(r)}$ such that
whenever $t,t'$ differ in one coordinate, the chosen cycles have at
least two common edges.
\end{lemma}
\begin{proof}
Supplementary Appendix~E specifies five and fifteen literal cycle templates
for $r=4$ and $r=5$, respectively. Normalize colours in first-appearance
order, append unused colours in increasing order, and restore each word
by the inverse palette map. The exact reader checks coverage and cyclic
adjacency for all 64 and 256 labelled boundaries, then intersects the chosen
cycle edge sets at all 576 and 3,072 directed one-coordinate interfaces.
The minimum intersections have three and five edges, respectively.
In particular, changing $a$ changes both endpoint exclusions in one step.
\end{proof}

\begin{lemma}\label{lem:shared-hub-paths}
Let $G$ have girth at least five. Suppose $v_0\cdots v_{r-1}$ is a path,
where $r\in\{4,5\}$, its endpoints have degree two and a common neighbour
$w$ outside the path, and its internal vertices have degree three.
Put $H=G-\{v_0,\ldots,v_{r-1}\}$. If $\R_4(H)$ is Hamiltonian,
then $\R_4(G)$ is Hamiltonian.
\end{lemma}
\begin{proof}
There are no additional internal edges. The third neighbours of the internal
vertices are distinct and different from $w$: a prohibited coincidence
would create a triangle or a $4$-cycle. The extension fibres are therefore
exactly the graphs in Lemma~\ref{lem:shared-hub-path-data}.

Cut an edge of an arbitrary bottom Hamilton cycle and process its fibres in
path order, giving each its certified cycle. An outside move either leaves
the boundary unchanged or changes one physical coordinate. The chosen
cycles share at least two internal edges. Both endpoints of a common edge
are legal in both fibres, so the two matching cross edges perform the actual
outside recolouring.

The current last fibre has lost at most its one incoming splice edge.
Choose a remaining common edge, remove it from the current cycle and the
fresh fibre cycle, and add the two matching cross edges. This joins the
cycles. Iteration gives one cycle through every extension of every bottom
colouring, and hence a Hamilton cycle of $\R_4(G)$.
\end{proof}

\subsection{Further degree-two reductions}

A shared-hub path has degree-two endpoints with the same neighbour outside
the path. For $m\in\{5,6\}$, a tailed $m$-cycle consists of an induced
cycle $v_0\cdots v_{m-1}v_0$ and an additional vertex $u$ whose only
neighbour on the cycle is $v_0$, where
\[
 d(v_0)=4,\qquad d(v_1)=d(u)=2,\qquad
 d(v_i)=3\quad(2\le i<m).
\]
The other neighbour of $u$ is outside this configuration. All degrees
in the following lemma are degrees in $G$.

\begin{lemma}\label{lem:further-degree-two}
Let $G$ have girth at least five, and let $X$ be one of the following
configurations:
\begin{enumerate}
\item a degree-four vertex and three of its degree-two neighbours;
\item an induced shared-hub path with degree sequence
$(2,4,3,2)$, $(2,4,3,3,2)$, $(2,3,4,3,2)$, or
$(2,3,3,3,3,2)$, allowing reversal;
\item a tailed $5$-cycle or tailed $6$-cycle.
\end{enumerate}
If $\R_4(G-X)$ is Hamiltonian, then $\R_4(G)$ is Hamiltonian.
\end{lemma}
\begin{proof}
Supplementary Appendix~F specifies seven literal cycle families for these
configurations. A port denotes an outside-neighbour role; for a shared-hub
path, one port controls both endpoints. Colours are normalized by first
appearance, with unused colours appended in increasing order. The inverse
palette map selects one fixed cycle at every labelled boundary.

Two distinct port roles can denote the same outside vertex only if all
pairs of their local neighbours have internal distance at least three;
otherwise a triangle or a $4$-cycle results. Appendix~F lists all allowed
pairs. Their compatibility graphs have no triangle, so no three distinct
roles can coincide. Several disjoint pairs may coincide, but an actual
outside move changes only one singleton or one listed pair. Both entries
of a pair change together from the same old colour to the same new colour.

The exact reader reconstructs the seven fixed internal graphs and all legal
states, checks the 421 normalized cycles, and verifies all 8,704 labelled
boundaries and 158,976 directed interfaces, including the paired changes.
The chosen cycles share at least two edges at every interface. A separate
reader restores every cycle into physical colour labels and directly checks
the same interface inequalities. Table~F.1 gives the individual counts.

Cut an edge of a bottom Hamilton cycle and process its fibres in path
order. Consecutive fibres have at least two common chosen-cycle edges;
an unchanged boundary has identical chosen cycles. The current last fibre
has lost at most its incoming splice edge. Choose another common edge,
remove its two copies, and add the two matching cross edges, which perform
the actual outside recolouring with the local state fixed. This joins two
cycles and preserves the invariant. Iteration gives one cycle through all
extensions of all bottom colourings, as required.
\end{proof}

\subsection{Linked pentagons and three hub contacts}

Two ordinary pentagons sharing exactly one edge have the model
\[
 (0,1,4,3,2,0),\qquad (0,1,7,6,5,0),
\]
possibly with the extra edge $36$. The shared endpoints already have
degree three; among the remaining vertices, only $36$ is compatible
with girth five. The exterior roles are $(2,3,4,5,6,7)$ in the bare
model and $(2,4,5,7)$ in the chorded model. Each role excludes one
outside colour; every other list is full.

We also use two disjoint ordinary pentagons joined by one edge, or by
two edges whose endpoints are nonadjacent on each pentagon. Their
models are the cycles $012340$ and $567895$, with edge $05$ and,
in the second case, edge $27$. Finally, the three-contact path has
vertices $v_0,\ldots,v_6$ and exclusion pattern
\[
 (a,b,c,a,d,e,a).
\]
Here $a$ is one physical outside role, with contacts at $v_0,v_3,v_6$.

\begin{lemma}\label{lem:extended-cycle-data}
The five models above have simultaneously selected Hamilton cycles
at every labelled boundary, sharing at least two edges across every
admissible outside recolouring.
\end{lemma}
\begin{proof}
Supplementary Appendix~H supplies the literal cycles. Its two readers
reconstruct the legal states, check exact cyclic coverage, and replay
every physical interface. The numbers of templates for the shared-edge
bare and chorded models, the three-contact path, and the one-link and
two-link models are respectively $187,15,51,451,67$. Their observed
minimum common-edge counts are $56,94,10,219,390$.

The first three use first-appearance palette normalization, with unused
colours appended increasingly. The linked models also minimize over
their fixed spatial symmetries. The inverse maps select one cycle at
each labelled boundary. An independent physical-label reader reproduces
the full interface histograms.

Two distinct exterior roles may coincide only when all pairs of their
contacts have internal distance at least three. The readers reconstruct
these compatibility graphs and check their complete alias matchings.
None has a compatible triple, so a move changes one role or one equal
pair simultaneously. This includes the single three-list hub move on
the path. Static identifications of other disjoint pairs do not add
another move type. All $75{,}008$ labelled boundaries and $2{,}577{,}024$
directed interfaces are covered; Appendix~H gives the individual counts.
\end{proof}

\begin{lemma}\label{lem:extended-cycle-restoration}
Let $G$ have girth at least five, and let $X$ be one of the following
induced configurations:
\begin{enumerate}
\item the bare or chorded shared-edge model, with all eight vertices
of degree three in $G$;
\item two disjoint ordinary pentagons with exactly one joining edge,
or exactly two joining edges with nonadjacent endpoints on each cycle;
\item a path $v_0\cdots v_6$ with degree-two endpoints, all other
vertices of degree three, and a vertex $w$ outside it adjacent to
$v_0,v_3,v_6$.
\end{enumerate}
If $\R_4(G-X)$ has a Hamilton path, then $\R_4(G)$ is Hamiltonian.
\end{lemma}
\begin{proof}
The actual extension fibres have exactly the prescribed models and
outside roles. In the last case girth excludes any further contact
of $w$ with the path; its other neighbours are unrestricted.
Give each fibre its selected cycle and process the bottom Hamilton
path in order. An unchanged boundary has identical cycles. Otherwise
Lemma~\ref{lem:extended-cycle-data} supplies at least two common edges.
The current last fibre has lost at most its incoming splice edge, so
another common edge remains. Remove its two copies and insert the two
matching cross edges, which perform the actual outside recolouring
with a local assignment legal on both sides. Each switch joins two
cycles. Iteration covers all extensions of all bottom colourings.
\end{proof}

\begin{lemma}\label{lem:any-two-ordinary}
Let $C,D$ be distinct ordinary facial pentagons in a plane graph $G$
of girth at least five. There is a set $X\subseteq V(C)\cup V(D)$
such that Hamiltonicity of $\R_4(G-X)$ implies Hamiltonicity of
$\R_4(G)$.
\end{lemma}
\begin{proof}
If the two faces intersect, every shared vertex lies on a shared edge,
since it has degree three. Shared edges cannot be consecutive, since
only one facial angle uses a given pair of edges at such a vertex.
If there are $r$ common edges, their union has $10-2r$ vertices and
$10-r$ edges. The planar girth-five edge bound gives
\[
 10-r\le\frac53(8-2r),\qquad 7r\le10,
\]
so $r=1$. Its induced union is the bare or chorded shared-edge model,
and Lemma~\ref{lem:extended-cycle-restoration} applies.

Otherwise the cycles are disjoint, and their joining edges form a
matching. If two matching endpoints are adjacent on one cycle, their
mates are nonadjacent on the other, by girth. The adjacent edge, the
two joining edges and a two-edge arc between those mates form an
ordinary pentagon sharing exactly one edge with the first cycle.
Its induced shared-edge union is again one of the two checked models;
the constructed pentagon need not be facial.

In the remaining case the matching endpoints are independent on both
cycles, so there are at most two joining edges. With no joining edge,
Lemma~\ref{lem:two-pentagons-restore} applies. With one or two, use
Lemma~\ref{lem:extended-cycle-restoration}.
\end{proof}

\subsection{Pentagonal reservoirs}

\begin{lemma}\label{lem:pentagonal-reservoir-data}
With palette $Q=\{0,1,2,3\}$, the following models admit fixed Hamilton
cycles in every component of every boundary fibre. The last column lists
the outside-neighbour roles whose colours each local vertex must avoid;
cycle vertices are ordered $v_0,\ldots,v_4$, followed by $z$ when present.
\[
\begin{array}{c|c|l}
\text{model}&\text{internal graph}&\text{forbidden role sets}\\ \hline
R&C_5\sqcup\{z\}&\{0\},\{1\},\{2\},\{3\},\{4\},\{5,6\}\\
F&C_5&\varnothing,\{0\},\{1\},\{2\},\{3\}\\
P&C_5+v_0z&\varnothing,\{0\},\{1\},\{2\},\{3\},\{4\}
\end{array}
\]
Every fibre of $F$ and $P$ is connected. A fibre of $R$ has two components
exactly when its first five boundary entries are equal, and is otherwise
connected. Across every admissible physical boundary move, each component
on either side overlaps a component on the other side. For every overlapping
pair, the selected component cycles have at least four common edges in
$R$, and at least two in $F$ and $P$.
\end{lemma}
\begin{proof}
Supplementary Appendix~G specifies the three literal cycle families.
Normalize colours by first appearance, appending unused colours in
increasing order, and restore words by the inverse palette map. This selects
the cycles simultaneously at all labelled boundaries. The exact readers
reconstruct the legal states and components, check cyclic coverage and
adjacency, and test all labelled interfaces. A separate physical-label
reader reproduces the interface histograms.

The three families have respectively $715,15,51$ normalized boundaries,
$720,15,51$ component cycles, and $16{,}384,256,1{,}024$ labelled boundaries.
The numbers of directed boundary interfaces are
$466{,}944,3{,}072,16{,}896$, with minimum common-edge counts $4,11,8$.
For $R$, all $470{,}208$ overlapping component interfaces are checked.

In $R$, the first five roles are distinct, as are the last two. The possible
extra identifications are the pairs $(i,j)$ with $0\le i<5$ and
$j\in\{5,6\}$. Their simultaneous identifications form the 31 matchings
of this bipartite graph. In $F$ all roles are distinct. In $P$ the only
possible pairs are $(1,4)$ and $(2,4)$, and they cannot both coincide.
Girth excludes every other identification. Checking all labelled single-role
and simultaneous equal-pair moves covers these possibilities, including
other static identifications in the same boundary.
\end{proof}

\begin{lemma}\label{lem:pentagonal-reservoirs}
Let $G$ have girth at least five, and let $X$ be one of the following:
\begin{enumerate}
\item an ordinary pentagon and a degree-two vertex outside it with no
neighbour on the pentagon;
\item a pentagon with one degree-two vertex and all other vertices of
degree three;
\item an ordinary pentagon and an adjacent degree-two vertex outside it.
\end{enumerate}
If $\R_4(G-X)$ is Hamiltonian, then $\R_4(G)$ is Hamiltonian.
\end{lemma}
\begin{proof}
Girth excludes chords and gives exactly the three models of
Lemma~\ref{lem:pentagonal-reservoir-data}. In the last case the outside
degree-two vertex meets the pentagon only once; deleting it leaves a full
list at that cycle vertex. Every matching cross edge at a common local
state performs the actual bottom recolouring, since that state is legal
in both fibres.

Cut a bottom Hamilton cycle into a path and give each component of each
fibre its fixed selected cycle. For the two connected-fibre models, splice
successive cycles in path order. The current last fibre has lost at most
one incoming splice edge. Of at least two common cycle edges with the next
fresh fibre, one remains; deleting its two copies and adding the matching
cross edges joins the cycles. An unchanged boundary uses identical cycles.

For $R$, form a graph whose nodes are the components of the layers, joining
overlapping components in consecutive layers. It has maximum degree at
most four. Some layer is connected. Otherwise the five cycle-neighbour
colours form a constant word at every bottom state. A physical move changes
at most one of these five entries, so that constant word is unchanged along
the bottom path. A palette permutation changing its colour produces another
proper bottom colouring, contradicting Hamilton coverage.

Every component can step to a component in either adjacent layer. Walking
towards a connected layer therefore proves that the component graph is
connected. Choose a rooted spanning tree and attach each fresh component
cycle to its already attached parent. Before an attachment the parent has
lost at most three internal edges, since its tree degree is at most four
and one attachment remains. At least four common selected-cycle edges are
available, so an uncut edge remains for the switch. An unchanged boundary
uses identical component cycles, which also have at least four edges.
Each switch joins two cycles. The final cycle covers every component of
every fibre, hence all proper colourings of $G$.
\end{proof}

\subsection{Degree-two vertices in a minimum counterexample}

In the next two lemmas, a minimum counterexample means a graph of minimum
order with non-Hamiltonian $\R_4(G)$ in a fixed class of nonempty planar
graphs of girth at least five closed under nonempty vertex deletion.
Thus every smaller nonempty induced graph has a bottom Hamilton cycle.

\begin{lemma}\label{lem:degree-two-counterexample}
A minimum counterexample is connected and has minimum degree at least two.
Let $S$ be its set of degree-two vertices. If $|S|\ge2$, its members have
a common heavy neighbour. If $S\ne\varnothing$, every $5$-cycle meets a
heavy vertex.
\end{lemma}
\begin{proof}
The singleton case gives $K_4$. Products of Hamiltonian graphs are
Hamiltonian, by alternating two distinct splice edges along a factor
Hamilton path, so $G$ is connected. Lemma~\ref{lem:single-lift} excludes
degree at most one, and Lemma~\ref{lem:adjacent-two-four} makes $S$
independent.

Suppose $|S|\ge2$. If no outside vertex is adjacent to all of $S$, every
outside vertex has at most $|S|-1$ neighbours in $S$. Put $H=G-S$ and
choose $u\in S$ with neighbours $x,y$. They are nonadjacent by girth.
Three-colour $H-\{x,y\}$ by Gr\"otzsch's theorem~\cite{Grotzsch59}, and
give both $x,y$ colour four. This leaves three colours available at $u$,
so Lemma~\ref{lem:independent-restoration} contradicts minimality.
There is consequently a common neighbour $w$ outside $S$. It cannot have
degree two, since $S$ is independent, and degree three would give
Lemma~\ref{lem:path232}; the three external neighbours of that path
are distinct by girth. Thus $w$ is heavy. Write $N(u_i)=\{w,x_i\}$; the $x_i$ are distinct
because a repetition creates a $4$-cycle.

If $S=\{z\}$, a pentagon avoiding heavy vertices either contains $z$
and has all other degrees three, or is ordinary. Pair an ordinary pentagon
with $z$, using its adjacency or nonadjacency to the cycle. All three
possibilities are excluded by Lemma~\ref{lem:pentagonal-reservoirs}.
If $|S|\ge2$, a cycle containing a member of $S$ must contain $w$.
A pentagon avoiding heavy vertices is therefore ordinary; pairing it with
any member of $S$ gives the same contradiction.
\end{proof}

For a plane minimum counterexample write
\[
 s=|S|,\qquad h=|\{v:d(v)\ge4\}|,\qquad
 D=\sum_{d(v)\ge4}(d(v)-3),\qquad
 Q=\sum_{|F|\ge7}(|F|-6),
\]
and let $p$ count pentagonal faces. Lemma~\ref{lem:five-face-identity-subclass}
gives
\begin{equation}\label{eq:degree-two-euler}
 p=12+Q+2D-2s.
\end{equation}
Pentagonal faces give distinct $5$-cycles unless $G=C_5$, already covered
by Theorem~\ref{thm:subcubic}.

\begin{lemma}\label{lem:heavy-vertex-budgets}
If a minimum counterexample has $s=1$, then
\[
 10+D+Q\le3h,\qquad h\ge5,\qquad p\ge20.
\]
If $s\ge2$, write
\[
 N(u_i)=\{w,x_i\},\qquad d(w)=s+t,
\]
and let $h'$ count the heavy vertices other than $w$, with
$D'=\sum_{v\ne w,\ d(v)\ge4}(d(v)-3)$. Then $t\ge1$ and
\begin{align}
 p&=6+Q+2t+2D',\label{eq:degree-two-hub}\\
 D'+6+Q&\le3h',\qquad D'\ge h'\ge3.\label{eq:degree-two-heavy-budget}
\end{align}
\end{lemma}
\begin{proof}
If $s=1$, every pentagonal face meets a heavy vertex by
Lemma~\ref{lem:degree-two-counterexample}. Counting its incidences gives
$p\le3h+D$. Combine this with \eqref{eq:degree-two-euler} and $D\ge h$
to obtain all three assertions.

For $s\ge2$, substituting $D=s+t-3+D'$ gives
\eqref{eq:degree-two-hub}. A pentagonal face having no heavy vertex other
than possibly $w$ must contain $w$. It cannot use two edges from $w$ to
$S$: its other two cycle vertices would give the regular shared-hub path
of order four in Lemma~\ref{lem:shared-hub-paths}. Thus it uses an edge
from $w$ to $Z=N(w)\setminus S$, so there are at most $2t$ such faces.
Counting all remaining pentagonal faces at other heavy vertices gives
$p\le2t+3h'+D'$. This proves \eqref{eq:degree-two-heavy-budget}, since
$D'\ge h'$.

If $t=0$, then $s=d(w)\ge4$. Among the distinct $x_i$, two are nonadjacent,
since otherwise three form a triangle. Put $J=G-(S\cup\{w\})$.
Three-colour $J$ with these two vertices removed by Gr\"otzsch's theorem,
then give both colour four. Induction on $J$ and
Lemma~\ref{lem:subdivided-star-four-subclass} contradict minimality.
Therefore $t\ge1$.
\end{proof}

For $s\ge2$, retain the nonnegative deficit
\begin{equation}\label{eq:degree-two-deficit}
 R=2t+3h'+D'-p=3h'-D'-6-Q\ge0.
\end{equation}
Let $p_0$ count pentagonal faces with no heavy vertex other than possibly
$w$, and let $a(F)$ count incidences of heavy vertices other than $w$ on a
face $F$. Summing these incidences gives the exact decomposition
\begin{equation}\label{eq:degree-two-deficit-parts}
 R=(2t-p_0)+\sum_{|F|\ne5}a(F)
       +\sum_{|F|=5,\ a(F)>0}(a(F)-1).
\end{equation}
Every term is nonnegative. Incidences count repeated occurrences on a
nonpentagonal facial walk.

\begin{lemma}\label{lem:short-facial-walks}
In a connected plane graph of girth at least five and minimum degree at
least two, a nonsimple facial walk has length at least ten. Put
$Q=\sum_{|F|\ge7}(|F|-6)$. If $Q\le1$, every face has a simple boundary of length at most seven, and distinct
angles at any vertex belong to distinct faces.
\end{lemma}
\begin{proof}
A facial walk is cyclically nonbacktracking. Splitting it at a repeated
vertex gives two nonempty closed subwalks without internal immediate
reversal. Each contains a simple cycle: a nontrivial closed walk in a tree
reverses at an internal leaf. Both subwalks therefore have length at least
five. If $Q\le1$, no face has length at least eight, so every boundary is
simple. A simple boundary visits a vertex at most once, proving the last
assertion.
\end{proof}

\subsection{Protected pentagon restoration}
\begin{theorem}\label{thm:subcubic}
Let $G$ be a nonempty finite simple planar graph of girth at least five
and maximum degree at most three. Then $\Rk(G)$ is Hamiltonian for
every $k\ge4$. Disconnected graphs are allowed. The uniform Hamiltonian threshold
of this class is exactly four, and its Hamilton-connected threshold lies
between four and five.
\end{theorem}

We prove a local restoration lemma using only Hamiltonicity of the
smaller recolouring graph, and then supply the planar structure.
Supplementary Appendices~C--D specify the cycles, seeds, pools and
interface records for 2,304 labelled pentagon boundaries and 313
standard-source one-entry interfaces. We first prove the protected-edge
invariant, then use it in an arbitrary-order restoration.

\paragraph{Local certificates.}
Put $Q=\{0,1,2,3\}$. A boundary $t$ has five entries in $Q$, with at most
one entry $\ast$. Its local graph $F_t$ consists of the proper $C_5$
colourings $x$ satisfying $x_i\ne t_i$; $\ast$ excludes no colour.
In the certificates $\ast$ is encoded by 4.
The eleven ordinary types have 590 states in total.
The constant type has two components, each a 15-cycle; all other
ordinary types are connected. Eight have the printed Hamilton cycles,
and the two remaining types are $00001$ and $00112$.
With a full list, the eleven types, after rotation, reflection and
colour permutation, are
\[
0000\ast,\ 0001\ast,\ 0010\ast,\ 0011\ast,\ 0012\ast,\
0101\ast,\ 0102\ast,\ 0110\ast,\ 0112\ast,\ 0120\ast,\ 0123\ast.
\]
Their printed Hamilton cycles have respectively
$54,60,78,62,71,88,85,74,74,83,82$ states, a total of 811.
Thus there are 23 component types, including the two constant components.

We record the four local facts used in the restoration as Claims 1--4.

\begin{claim}\label{lem:pentagon-path-blocks}
Every one of the 23 component types $B$ has a partition into paths
$P_1,\ldots,P_r$ of order two or three, with an actual $B$ edge
joining $P_i$ to $P_{i+1}$ for each $i<r$.
\end{claim}
\begin{proof}
Twenty-one types have the printed Hamilton cycles, and $00112$ has
the 51-state Hamilton path in supplementary Appendix~C.
Delete a cycle edge where necessary, and divide the resulting path
into consecutive two-vertex blocks, with one three-vertex block if
its order is odd. For $00001$, supplementary Appendix~C supplies the ordered path blocks and their actual connectors.  The verifier checks coverage and every one-coordinate difference.
\end{proof}

For a Hamilton component $A$ with specified cycle $C_A$, designate an
ordered pool of five edges for either constant component, and eight
edges otherwise. Let $K_A$ be the other cycle edges. The first two
pool edges join copies of $C_A$ in a two- or three-row path product;
the remaining pool edges supply at least two splice reserves, all
outside $K_A$.
For a hard component, specify cycles of $A\square P_2$ and
$A\square P_3$, and a set $D_A$ of internal edges occurring in
\emph{every} row of both seeds. For each of $00001$ and $00112$,
supplementary Appendices~C--D specify a two-row cycle and derive its three-row
seed by the same three-path insertion rule. Both use $|D_A|=13$
and a two-edge pool. Put $K_A=D_A\setminus\text{pool}$; the two
distinct pool edges supply the splice reserves.

Here is the exact normalization convention. Enumerate $s=0,\ldots,4$,
then $\epsilon=+1,-1$, then the permutations $p\in S_4$ in lexicographic
order. Take the first standard type
$n_i=p(t_{s+\epsilon i})$, retaining $\ast$. Restoration puts the
standard entry $x_i$ at position $s+\epsilon i$ with colour $p^{-1}(x_i)$.
In particular a standard boundary first normalizes by the identity.

\begin{claim}\label{lem:pentagon-protection}
Let the source be a standard component $A$ of $F_t$, and change one
non-star boundary entry. Normalize the target by the preceding rule.
Every target component $B$ whose restored state set meets $A$ has an
edge in common with $K_A$ after restoration of $K_B$.
For an unchanged boundary use the same component and the same map.
\end{claim}
\begin{proof}
The ordinary and starred orbit representatives cover
$4^5+5\cdot4^4=2304$ labelled boundaries. The displayed cyclic words
give their specified component cycles and the two-row hard seeds:
their entries are precisely the legal words, without repetition,
and successive entries, including last and first, differ in one
coordinate. Both three-row seeds follow from the supplement's
three-path insertion rule, using two edges and the displayed
40-state or 47-state path. Every listed common edge occurs in
every row of both corresponding seeds.
Removing the listed pools therefore defines the asserted protected sets.
For the 313 standard-source interfaces, supplementary Appendix~D gives 47
individual common edges, twenty-one source-selector rows covering 236
cases, and six constant-source families covering 30 cases.
Endpoint membership in the specified templates and outside their pools
proves each individual record. The selector rows use fixed protected
incident edges against the opposite component's full degree, with the
prescribed first normalizer and pools. The constant-source families use
two normalizing directions and four protected target edges. This is a
complete finite case proof for standard sources; labelled sources are
obtained by transport, with the target normalized relative to the source
as described after Lemma~\ref{lem:pentagon-overlap}.
\end{proof}

\paragraph{Products and overlaps.}
\begin{lemma}\label{lem:protected-product}
For any two of the 23 component types $A,B$, the graph $A\square B$
has a Hamilton cycle retaining a copy of every edge of $K_A$
in every physical $B$ row.
\end{lemma}
\begin{proof}
Use the ordered path blocks $P_1,\ldots,P_r$ of
Claim~\ref{lem:pentagon-path-blocks} for $B$.  For a two-row block
with rows $b_0,b_1$, remove the same pool edge $xy$ from the two
copies of $C_A$ and add
\[
 (x,b_0)(x,b_1),\qquad (y,b_0)(y,b_1).
\]
This is one cycle on $A\square P_2$.  For a three-row block use two
distinct pool edges $x_1y_1,x_2y_2$ and perform this operation between
rows $b_0,b_1$ and $b_1,b_2$, respectively.  The two cuts in the
middle row are distinct, so the result is again one cycle.  For either
hard component, use the literal two- or three-row seed; the
certificate checks that every edge of $D_A$ occurs in every physical
row of that seed.  In all cases every edge of $K_A$ survives in every
row.  For an ordinary component retain two of the unused pool edges
as reserves; for a hard component the two pool edges in its seed are
the reserves.  Thus every reserve is outside $K_A$.

Fix two such reserves $e$ and $f$.  Between $P_i$ and $P_{i+1}$,
let $u_i\in P_i$ and $v_i\in P_{i+1}$ be the certified adjacent
states.  Cut the copies of $e$ in the two rows $u_i,v_i$ when $i$ is
even, and the copies of $f$ when $i$ is odd; add the two vertical
edges obtained by holding the $A$-coordinate fixed and changing
$u_i$ to $v_i$.  The selected edge in each block is then different
at its two interfaces whenever both interfaces use the same physical
row.  Thus every cut is present, no edge of $K_A$ is cut, and each
operation joins two cycles into one.  Induction on $i$ gives a single
cycle on all rows, including the one-block case.  All added edges are
actual one-coordinate changes in $B$.
\end{proof}

\begin{lemma}\label{lem:pentagon-overlap}
For a fixed star position, or no star, let two boundaries be equal or
differ in one non-star position. Any overlapping component pair has
at least ten common vertices.
\end{lemma}
\begin{proof}
An unchanged constant component has 15 vertices. Between a constant
and nonconstant boundary, one additional colour is excluded at one
coordinate of that 15-cycle. A cyclic permutation of its three colours
preserves each component: an order-three action on a two-element
component set is trivial. Hence each colour appears five times at
each coordinate and precisely ten states remain.
In every other case, both components are their entire local state
sets. The intersection lists have size at least three except possibly
one list of size two. Start at that position, proceed around the
cycle, and finally avoid the two already coloured neighbours.
This gives at least $2\cdot2^3\cdot1=16$ colourings. With no changed
position the bound only improves. Full lists cause no difficulty.
\end{proof}

\paragraph{Relative restoration maps.}
Call a coordinate and colour map $\Psi$ a restoration map for a
boundary $t$ when $\Psi^{-1}(t)$ is standard; it restores the
specified cycles, pools and protected sets of that standard type. Let
$t'$ equal $t$ or differ from it in one non-star entry. Put
$\Psi'=\Psi$ if $t'=t$, and otherwise $\Psi'=\Psi\circ N_u$, where
$u=\Psi^{-1}(t')$ and $N_u$ is the canonical restoration map of $u$.
Transporting Claim~\ref{lem:pentagon-protection} by $\Psi$ shows that
overlapping components at $t$ and $t'$, restored by $\Psi$ and
$\Psi'$, have a common protected edge; the verifier also replays this
rule over all labelled boundaries and admissible maps $\Psi$. The
dependence of $\Psi'$ on $\Psi$ is necessary: for the labelled
boundaries $00002$ and $00003$, both of type $00001$, the two canonical
maps restore protected sets with no common edge, although the
components overlap.

Now let two consecutive layers carry maps related in this way in each
factor; the two factors use independent maps and may change at the same
base edge. If $xx'$ is a common protected edge of overlapping
first-factor components and $y_1,\ldots,y_{10}$ are common states of
overlapping second-factor components, then every
$\{(x,y_i),(x',y_i)\}$ is an internal edge of both product cycles of
Lemma~\ref{lem:protected-product} with the first factor as cycle
factor, since these retain every protected edge in every row. Thus the
two product cycles share at least ten internal edges.

\begin{claim}\label{lem:pentagon-component-graph}
Let $h_0h_1\cdots h_{m-1}$ be a Hamilton path in the
recolouring graph of the graph outside two pentagons.  For
$j\in\{1,2\}$, let $T_j(i)$ be the boundary list seen by the $j$th
pentagon at $h_i$, and let $c_j(i)=1$ when $F_{T_j(i)}$ is a constant
two-component fibre, and $c_j(i)=0$ otherwise.  Form the graph whose
vertices are the product components of
$F_{T_1(i)}\square F_{T_2(i)}$, $0\le i<m$, and whose edges record a
common local state across $h_ih_{i+1}$, $0\le i<m-1$.  This component
graph is connected.
\end{claim}
\begin{proof}
Give the two components of a constant fibre labels $0,1$ and give a
nonconstant fibre the single label $\ast$.  A base edge changes at
most one entry in either boundary list.  If both consecutive fibres of
a factor are constant, their boundary lists are equal and the label is
preserved; otherwise, by Lemma~\ref{lem:pentagon-overlap}, every label
on one side meets every label on the other.  The two factors follow
these rules independently, and this describes the product component
graph exactly.

For each $j$, some $r_j$ has $c_j(r_j)=0$.  Otherwise every $T_j(i)$
is a constant word.  Two constant words differing in at most one
position are equal, so the whole Hamilton path would have one fixed
boundary word.  A global permutation of the four colours changing that
word would give another proper colouring of the outside graph, which
the Hamilton path must contain, a contradiction.

From any product component, walk layer by layer to $r_1$ and then to
$r_2$, choosing label $0$ whenever a constant fibre is entered from a
nonconstant one.  The first label is $\ast$ at $r_1$ and then evolves
by this fixed rule, while the second label is $\ast$ at $r_2$.  Hence
every walk ends at the same component of layer $r_2$, and the graph is
connected.
\end{proof}

\begin{claim}\label{lem:protected-tree-splicing}
Let $Q$ be a connected graph of maximum degree at most $\Delta$.  For
each $v\in V(Q)$ let $X_v$ be a vertex-disjoint fibre carrying a
Hamilton cycle $C_v$.  Suppose that, for every $uv\in E(Q)$, the two
reference cycles share at least $r\ge\Delta$ internal edges, and
each shared edge has the two corresponding cross edges in the ambient
graph.  Then the union of the fibres has a Hamilton cycle.
\end{claim}
\begin{proof}
Root a spanning tree of $Q$ and process its vertices from the root
outwards.  Maintain one cycle on the fibres already processed; the
only edges removed from a processed fibre are the internal edges used
for its earlier tree attachments.  When a child $v$ is attached to
its parent $u$, at most $\deg_Q(u)-1\le\Delta-1$ shared edges at $u$
have been removed.  Since $r\ge\Delta$, a shared edge $xy$ remains
at $u$ and is still present in the fresh cycle $C_v$.  Delete $xy$
from both cycles and add the two cross edges at $x$ and $y$.  This
joins the two cycles into one and preserves the invariant.  Each
tree edge is processed once, so the final cycle covers every fibre.
\end{proof}

\paragraph{Restoration over a base cycle.}
\begin{lemma}\label{lem:two-pentagons-restore}
Suppose $G$ has girth at least five and contains two vertex-disjoint
5-cycles with no edge between them. Each cycle has vertices of degree
three except for at most one vertex of degree two.
Let $H$ be obtained by deleting both cycles. If $\R_4(H)$ is
Hamiltonian, so is $\R_4(G)$.
Neither planarity nor three-colourability of $H$ is assumed here.
\end{lemma}
\begin{proof}
Each degree-three cycle vertex has one exterior neighbour; a
degree-two vertex gives a full list. An exterior vertex meets a
given 5-cycle at most once, since two attachments would create a
3- or 4-cycle. Thus a bottom recolouring changes at most one list
in each factor; the two factors may change simultaneously.
All exterior neighbours remain in $H$, which is nonempty.

Delete one edge of a bottom Hamilton cycle to obtain a Hamilton path
$h_0,\ldots,h_{m-1}$. In each factor give $T_j(0)$ its canonical
restoration map, and propagate maps along the path by the relative
rule above. For every product component at $h_i$, take the product
cycle of Lemma~\ref{lem:protected-product} restored by the two maps
at $h_i$, with the first pentagon as cycle factor. The component graph
over the path is connected by Claim~\ref{lem:pentagon-component-graph}.
It has at most four vertices per layer, hence maximum degree at most
eight. Along each of its edges the two product cycles share at least
ten internal edges. A cross edge at either endpoint of such an edge is
the actual bottom recolouring, because the complete local colouring is
held fixed while one exterior vertex changes colour.
Claim~\ref{lem:protected-tree-splicing} with $\Delta=8$ and $r=10$
gives one cycle through all extensions of all bottom colourings.
\end{proof}


\subsection{Protected bowtie products}
A bowtie consists of two pentagons meeting only at a vertex of original degree four,
with all eight other vertices of original degree three. All degrees in
these definitions are measured before deletion.
Label its cycles $(0,1,2,3,4,0)$ and $(0,5,6,7,8,0)$. Girth excludes
chords within either cycle. An extra edge can only join a vertex of
$\{2,3\}$ to one of $\{6,7\}$. Two such edges either exceed a prescribed
degree or form a $4$-cycle with edges $23$ and $67$. Up to reflecting the
arms, the induced union is therefore the bare bowtie or the same graph
with edge $26$ added. In the bare model the exterior roles are
$1,\ldots,8$; in the chorded model they are $1,3,4,5,7,8$. Each role
excludes its outside colour. The other local lists are full, including
those at vertices $2,6$ of the chorded model.

\begin{lemma}\label{lem:clique-block-product}
Let $A$ be a graph, let $K\subseteq E(A)$, and let
$e,f\in E(A)\setminus K$ be distinct. Suppose $A\square K_2$ and
$A\square K_3$ have specified Hamilton cycles retaining
$K\cup\{e,f\}$ in every physical row. Suppose $B$ has a partition
into cliques of orders two and three, with actual $B$-edges joining
the blocks as a tree and at most two connector incidences at each
physical $B$-vertex. Then $A\square B$ has a Hamilton cycle retaining
every edge of $K$ in every physical $B$-row.
\end{lemma}
\begin{proof}
The connector graph on the actual vertices of $B$ is a forest: a cycle
would project to a closed trail with distinct edges in the quotient tree.
Its maximum degree is two, so label the edges of each nontrivial component
alternately $e,f$. Put the appropriate two- or three-row seed on each
clique block. For a connector $xy$ labelled $e=ab$, remove the copies
$(a,x)(b,x)$ and $(a,y)(b,y)$, and add $(a,x)(a,y)$ and $(b,x)(b,y)$.
Use $f$ for the other label. No row uses the same reserve twice, so each
cut remains present. The quotient tree ensures that each switch joins
two distinct current cycles. No edge of $K$ is removed, and the final
cycle covers $A\square B$. Neither $K$ nor the two reserves need form
a matching.
\end{proof}

\begin{lemma}\label{lem:bowtie-product-data}
For every boundary $t$ of either bowtie model, its local recolouring
graph $B_t$ has the following simultaneous finite data.
\begin{enumerate}
\item Hamilton cycles of $B_t\square K_2$ and $B_t\square K_3$
retain a matching $M_t$ in every physical row. Two specified edges
$e_t,f_t\in M_t$ are reserves, and $K_t=M_t\setminus\{e_t,f_t\}$
is the protected core.
\item The vertices of $B_t$ have an edge/triangle partition with a
connector tree of physical connector degree at most two.
\item Across every actual outside move, the selected physical cores
have at least nine common edges in the bare model and twenty-four
in the chorded model. An unchanged boundary uses the same core.
\end{enumerate}
In particular every $B_t$ is connected, and consecutive bowtie fibres
have at least eighteen common states.
\end{lemma}
\begin{proof}
Supplementary Appendix~I supplies the two seed families and their
clique trees for all $451$ bare and $107$ chorded canonical boundaries.
Normalize colours by first appearance, append unused colours increasingly,
and choose the first minimizing spatial map among the eight bare or two
chorded symmetries. Restoring by this fixed map selects the physical
seeds, matching and tree at every labelled boundary.

The exact readers independently reconstruct all legal state sets, verify
each seed's complete coverage and closing edge, and check the matching
in every physical row of both seeds. The tree reader checks a partition
without repetitions, every pairwise clique edge, all connectors, the
quotient tree and the connector degree at each physical vertex.
In particular, a three-vertex path is insufficient for a triangle block.

An outside vertex meets each pentagon at most once. Its possible roles
within a bowtie are thus singletons or pairs at internal distance at
least three; no triple is possible. The readers check every such move
on all $4^8$ bare and $4^6$ chorded labelled boundaries, including the
simultaneous change of equal-coloured paired roles. There are
$2162688$ and $89088$ directed interfaces. After the two reserves are
removed, primary and independent physical-label checks give identical
interface histograms and the stated core margins. Unchanged boundaries
retain the same core, of at least the corresponding minimum size.
Finally, nine common matching edges have eighteen distinct endpoints.
\end{proof}

The ordinary-pentagon component seeds already used in
Claim~\ref{lem:pentagon-protection} retain $K_A$ and two distinct
reserves outside $K_A$ in every row of $A\square P_2$ and
$A\square P_3$. They are also valid seeds for the corresponding
clique products, since $P_3\subseteq K_3$. The ordinary seed extract
is rechecked in supplementary Appendix~I. Its protected sets are
arbitrary edge sets, as permitted by Lemma~\ref{lem:clique-block-product}.

\begin{theorem}\label{thm:pentagon-bowtie-restoration}
Let $G$ have girth at least five. Suppose an ordinary pentagon $C$
and a bowtie $X$ are vertex-disjoint, with no edge between them.
Put $H=G-(V(C)\cup V(X))$. If $\R_4(H)$ is Hamiltonian, then
$\R_4(G)$ is Hamiltonian.
\end{theorem}
\begin{proof}
Cut a bottom Hamilton cycle into a path. The fibre over each bottom
colouring is $A_t\square B_s$, where $A_t$ is the ordinary-pentagon
fibre and $B_s$ is the appropriate bowtie fibre. Propagate the ordinary
restoration maps along this path by the relative rule following
Lemma~\ref{lem:pentagon-overlap}; retain the map at an unchanged boundary.
For each component $A$ of $A_t$, apply
Lemma~\ref{lem:clique-block-product} with its ordinary seeds and the
certified clique tree of $B_s$. This gives a cycle in $A\square B_s$
retaining every $K_A$-edge in every physical bowtie row.

Consecutive overlapping $A$-components have a common protected edge by
Claim~\ref{lem:pentagon-protection}. The bowtie fibres have at least
eighteen common states by Lemma~\ref{lem:bowtie-product-data}. Their
selected product cycles therefore share at least eighteen internal edges.
Each has its two actual matching cross edges: the whole local assignment
is legal on both sides and is held fixed during the bottom recolouring.
This also covers an outside vertex meeting one vertex of $C$ and two
vertices of $X$. The two deleted factors have no joining edge, so both
fibres and interface intersections factorize despite shared outside
neighbours.

Form the component graph over this bottom path. Each layer has at most
two components, since $B_s$ is connected, so its maximum degree is four.
Every component crosses each adjacent interface, by the ordinary-pentagon
overlap property and the bowtie state overlap. Some ordinary fibre is
connected. Otherwise all five exterior colours of $C$ would be constant
at every bottom state. A move changes at most one of those entries, so
the constant word could never change. A global palette permutation
changing its colour produces another proper bottom colouring, contradicting
coverage by the Hamilton path. These five distinct exterior neighbours
belong to $H$, since $C$ and $X$ are nonadjacent, so $H$ is nonempty.

Walking toward that connected layer proves the component graph connected.
Claim~\ref{lem:protected-tree-splicing}, with $\Delta=4$ and $r=18$,
joins all its reference cycles. They cover precisely all proper
colourings of $G$.
\end{proof}

\begin{theorem}\label{thm:paired-bowties}
Let $G$ have girth at least five and contain two vertex-disjoint bowties
with no edge between them. Delete their eighteen vertices to obtain $H$.
If $\R_4(H)$ has a Hamilton path, then $\R_4(G)$ is Hamiltonian.
\end{theorem}
\begin{proof}
Use the first bowtie as factor $A$ and the second as factor $B$.
Lemmas~\ref{lem:clique-block-product} and \ref{lem:bowtie-product-data}
give a Hamilton cycle in every fibre $A_t\square B_s$, retaining the
first factor's core in every second-factor row. Across a bottom move,
at least nine first-factor core edges and eighteen second-factor states
are common. The two reference cycles therefore share at least
$9\cdot18=162$ internal edges. This remains true if one physical
outside vertex changes two roles in each bowtie: local assignments and
their intersections still factorize.

Process the bottom Hamilton path in order. The most recently attached
fibre has lost only its incoming splice edge. Choose another common
internal edge, remove its two copies and insert the two matching cross
edges. Both cross edges are genuine outside recolourings. Each switch
joins two cycles, and the final cycle covers all extensions of all bottom
colourings. The outside graph is nonempty: each bowtie has exterior
contacts and the other bowtie is nonadjacent. No bound on the bottom
path length is used.
\end{proof}

\begin{remark}
The bare boundary $(3,0,0,2,3,0,0,1)$ has $1432$ states and three leaves,
so its fibre has no Hamilton path. The new constructions use clique trees
and protected products; neither restoration assumes a spanning path in
a bowtie fibre.
\end{remark}

\begin{lemma}\label{lem:bowtie-degree-two}
Let $G$ have girth at least five. Suppose a bowtie $X$ and a degree-two
vertex $z$ outside it are disjoint and have no edge between them. If
$\R_4(G-(V(X)\cup\{z\}))$ has a Hamilton path, then $\R_4(G)$ is
Hamiltonian.
\end{lemma}
\begin{proof}
A fibre is $B_t\square K_A$, where $A$ is the available set at $z$,
of size two or three. Use the certified clique seed, identifying its
rows with the physical colours in $A$. It retains every edge of $M_t$
in every row. Across an outside move, $A\cap A'\ne\varnothing$:
the two neighbours of $z$ are distinct, and only one can change.
The bowtie matchings have at least nine common edges by
Lemma~\ref{lem:bowtie-product-data}, even after its reserves are
removed. In any common colour row these give at least nine common
seed edges. This remains valid when the move changes two bowtie
roles and one neighbour colour of $z$ simultaneously. If the bowtie
boundary is unchanged, the same matching is used.

Process the bottom path as in
Lemma~\ref{lem:extended-cycle-restoration}. Only one incoming edge
has been cut at the newest fibre, so a common edge remains for each
next switch. The result is one cycle through all extensions. No
claim that an isolated bowtie fibre has a Hamilton path is needed.
\end{proof}


\subsection{Pentagons sharing an edge at a heavy vertex}
A \emph{hinged pair} consists of the pentagons
\[
 (0,2,3,4,1,0),\qquad (0,1,7,6,5,0),
\]
with vertex $0$ of original degree four and all seven other vertices
of original degree three. Its induced graph may also have edge $36$.
The exterior roles are $(0,2,3,4,5,6,7)$ in the bare model and
$(0,2,4,5,7)$ in the chorded model. Each role excludes one outside
colour. Vertex $1$, and the chord endpoints when present, have full lists.

Two facial pentagons with the same sole degree-four heavy vertex
form a hinged pair when their angles there are adjacent. At a common
light vertex the faces share an edge, and common edges cannot be
consecutive. If they share $r$ edges, their union has $10-2r$ vertices
and $10-r$ edges. The planar girth-five bound gives
\[
 10-r\le\frac53(8-2r),\qquad 7r\le10,
\]
so $r=1$. Among unsaturated local vertices only $36$ can be added
without creating a triangle or a $4$-cycle. Opposite facial angles
instead give the bare or chorded bowtie defined above: another common
vertex would give two distinct paths of length at most two from the
centre, and hence a triangle or $4$-cycle.

\Needspace{9\baselineskip}
\begin{lemma}\label{lem:hinged-data}
Every labelled boundary $t$ of either hinged-pair model has a selected
matching $M_t$ and Hamilton cycles of $A_t\square K_2$ and
$A_t\square K_3$ retaining $M_t$ in every physical row, where $A_t$
is the local recolouring fibre. Two specified edges of $M_t$ are
reserves; the remaining core $K_t$ has at least six common edges
across every bare interface and twenty across every chorded interface.
Every $A_t$ also has an edge/triangle partition with a connector tree
of physical connector degree at most two. Thus the fibres are connected,
and consecutive fibres have at least twelve common states.
\end{lemma}
\begin{proof}
Supplementary Appendix~J supplies unchanged literal seeds and trees.
The bare and chorded models have respectively $379$ and $31$ canonical
boundaries, with $241774$ and $26051$ states in total. Normalize by
first appearance after taking the minimum over the identity and the
arm exchange $(0,1,5,6,7,2,3,4)$; append unused colours increasingly.
Minimizing the full normalization tuple fixes all ties. Its inverse
selects physical data at every labelled boundary.

The two readers independently reconstruct all legal states and check
complete seed coverage, legality, closure and matching retention in
every row. The primary reader also checks every clique partition,
connector, quotient tree and physical connector capacity. An ordered
three-vertex path does not suffice for a triangle block.

An outside vertex can occupy only one role or two whose contacts
have internal distance at least three. The complete bare role pairs are
\[
 (1,5),(1,6),(2,4),(2,5),(2,6),(3,4),(3,5),
\]
and the chorded pairs are $(1,4),(2,3)$. The compatibility graphs
have no triangle; their complete alias matchings number $22$ and $4$.
Both readers check all $16384$ and $1024$ labelled boundaries and
all $430080$ and $16896$ directed interfaces, including simultaneous
changes at paired roles. The independent reader restores literal
physical strings at every boundary and reproduces the full matching
and post-reserve core histograms. Their core minima are six and
twenty; unchanged boundaries use identical cores of sufficient size.
Six common matching edges give twelve distinct common states.
\end{proof}

\begin{lemma}\label{lem:hinged-restorations}
Let $G$ have girth at least five and contain a hinged pair $X$.
Suppose $X$ is separated from a degree-two vertex, another hinged
pair, a bowtie, or an ordinary pentagon, where separated means
vertex-disjoint with no edge between the sets. Delete both displayed
configurations to obtain $H$. If $\R_4(H)$ has a Hamilton path,
then $\R_4(G)$ is Hamiltonian. Permitted chords and shared outside
neighbours are allowed.
\end{lemma}
\begin{proof}
For a degree-two vertex $z$, the fibre is $A_t\square K_L$, where
$L$ is its available set of two or three physical colours. Use the
corresponding clique seed. An outside move leaves $L\cap L'\ne
\varnothing$, since the neighbours of $z$ are distinct and only one
can change. In a common row the seeds retain at least six common
matching edges, even if the same move changes two roles of $X$.
The bottom contains the two neighbours of $z$, so is nonempty.

For another hinged pair, apply Lemma~\ref{lem:clique-block-product}
with $X$ as seed factor and the second factor's certified clique tree.
The resulting cycles retain the first core in every second-factor row.
Consecutive cycles share at least $6\cdot12=72$ internal edges.
For a bowtie use its clique tree and its eighteen common states from
Lemma~\ref{lem:bowtie-product-data}, giving at least $6\cdot18=108$
common internal edges. Fibre intersections factorize because the two
deleted factors have no joining edge. These counts therefore include
one outside move affecting two roles in each factor simultaneously.

In these three cases process the bottom Hamilton path in order.
The newest fibre has lost at most its incoming splice edge. Choose
another common internal edge, remove its copies and add the two
matching cross edges. Each cross edge performs the actual outside
recolouring with a fixed local assignment legal on both sides. Each
switch joins two cycles, and iteration covers all extensions.

For an ordinary pentagon $C$, use each ordinary component $B$ as the
seed factor and $A_t$ as the clique-tree factor. The ordinary
$P_2/P_3$ seeds are valid $K_2/K_3$ seeds. Propagate the ordinary
restoration maps along the bottom path by the relative rule following
Lemma~\ref{lem:pentagon-overlap}, retaining the map at an unchanged
boundary. Lemma~\ref{lem:clique-block-product} gives a product cycle
retaining every ordinary protected edge in every hinged row.
Overlapping ordinary components have a common protected edge, and
consecutive hinged fibres have twelve common states. Their product
cycles thus share at least twelve internal edges, including combined
two-role/one-role outside moves.

The component graph over the bottom path has at most two nodes per
layer and maximum degree at most four. Every component crosses each
adjacent interface. Some ordinary fibre is connected: otherwise its
five exterior colours would be a constant word at every bottom
colouring. A move changes at most one entry, so the word could not
change, contrary to a global palette permutation and Hamilton coverage.
Those five exterior neighbours belong to $H$. Walking to a connected
layer proves the component graph connected. Apply
Claim~\ref{lem:protected-tree-splicing} with $r=12$ and $\Delta=4$.
The resulting cycle covers all proper colourings of $G$.
\end{proof}

\begin{remark}
The bare boundary $(2,2,1,1,3,1,1)$ has $526$ states and three leaves,
so its fibre has no Hamilton path. Appendix~J checks this obstruction
as well as the protected product seeds; the restoration uses the products.
\end{remark}

\subsection{Sharper degree-two budgets}
Retain the minimum-counterexample convention and the notation of
Lemma~\ref{lem:heavy-vertex-budgets}. All degrees and facial incidences
below are measured in the original plane graph, before any deletion.

\begin{lemma}\label{lem:unique-degree-two-heavy-budget}
If a minimum counterexample has exactly one degree-two vertex, then it
has at least eight heavy vertices, $D\ge16$, and at least forty-two
distinct $5$-cycles.
\end{lemma}
\begin{proof}
Let $z$ have neighbours $x,y$. Every pentagon meets a heavy vertex by
Lemma~\ref{lem:degree-two-counterexample}, so
\[
 p=10+2D+Q,\qquad \Delta=3h+D-p=3h-D-10-Q\ge0.
\]
The deficit $\Delta$ counts nonpentagonal heavy incidences and excess
heavy incidences on pentagons. Thus at most $\Delta$ pentagons have
two or more heavy vertices.

Call a face near if it meets $N[z]$, and far otherwise. At most
$d(x)+d(y)-2$ pentagonal faces are near. If $zx,zy$ are not bridges,
their two distinct incident faces both contain $x,y$. If they are
bridges, their common bridge face repeats both $x$ and $y$, giving
the stronger bound $d(x)+d(y)-3$. Every face meeting $z$ meets $x,y$.
Put
\[
 B=\{x,y\}\cap\{v:d(v)\ge4\},\qquad
 D_B=\sum_{v\in B}(d(v)-3),
\]
and let $b_B$ count the vertices of $B$ of degree at least five.
The near-face bound is $4+D_B$.

Two sole-heavy far pentagons at a degree-four vertex form a hinged
pair or a bare/chorded bowtie, according as their angles are adjacent
or opposite. They are separated from $z$, contradicting
Lemma~\ref{lem:hinged-restorations} or
Lemma~\ref{lem:bowtie-degree-two}. Hence the sole-heavy far capacity
is one at degree four and at most $d(v)$ at higher degrees. If $a$
counts degree-four vertices, the total capacity is
\[
 C=a+\sum_{d(v)\ge5}d(v)=3h+D-3a.
\]
Let $C_B$ be its contribution at $B$. Since $D_B-C_B=-3b_B$,
\[
 p\le4+D_B+C-C_B+\Delta
 \quad\Longrightarrow\quad
 2D+3a+16+2Q+3b_B\le6h.
\]
Write $e=\sum_{d(v)\ge6}(d(v)-5)$, so $D=2h-a+e$. The last
inequality gives
\[
 D\ge16+2Q+3e+3b_B,\qquad
 2h\ge16+a+2e+2Q+3b_B.
\]
Thus $D\ge16$, $h\ge8$, and $p\ge42$. These pentagonal faces give
distinct $5$-cycles, as noted after \eqref{eq:degree-two-euler}.
\end{proof}

Suppose now that $s\ge2$. Write $N(u_i)=\{w,x_i\}$ and
$Z=N(w)\setminus S$, where $|Z|=t\ge1$. Recall that $p_0$ counts
pentagonal faces having no heavy vertex other than possibly $w$.
Such a face contains $w$ and cannot use an $S$--$S$ angle there.
Set $b=2t-p_0$.

\begin{lemma}\label{lem:degree-two-z-runs}
Let the cyclic order of the neighbours of $w$ be written as an
$S/Z$ word. If a $Z$-vertex is flanked by $S$-vertices, its two mixed
faces cannot both be counted by $p_0$. If the lengths of the maximal
$Z$-runs are $\ell_1,\ldots,\ell_r$, then
\begin{equation}\label{eq:degree-two-z-run-cost}
 b\ge\sum_{j=1}^r\max\{1,\ell_j-1\}\ge1.
\end{equation}
\end{lemma}
\begin{proof}
Suppose the two faces beside $wz$, with $z\in Z$ isolated in the word,
are pentagons
\[
 (w,u,x,a,z,w),\qquad (w,z,b,y,v,w),
\]
where $u,v\in S$ are distinct. Every vertex other than $w,u,v$
displayed here has degree three. A further degree-two vertex on either
pentagon would have to use its edge to $w$ on that cycle, contrary to
the mixed angle.

The seven vertices $u,x,a,z,b,y,v$ are distinct. An internal vertex
cannot coincide with $u,z,v$ without creating a triangle. The
equalities $x=y$, $x=b$, and $a=y$ create $4$-cycles through $w$.
Finally $a=b$ is impossible, since the two faces beside $wz$ use
different other neighbours of the degree-three vertex $z$. The only
possible chord of this seven-vertex path is $xy$: the degrees at
$u,v,z$ are exhausted, and any other edge among $x,a,b,y$ creates a
triangle or a $4$-cycle. If $xy$ exists, $u,x,y,v$ is the regular
shared-hub path of order four. Otherwise the induced path has the
three-contact outside neighbour $w$ at its first, fourth, and seventh
vertices. Lemma~\ref{lem:shared-hub-paths} or
Lemma~\ref{lem:extended-cycle-restoration}, respectively, contradicts
minimality.

There are $2t$ edge-side incidences at the $Z$-edges of $w$. Each
$p_0$-face uses one at a mixed angle or two at a $Z$--$Z$ angle. Thus
each $Z$--$Z$ angle costs at least one in $b$, and an isolated
$Z$-vertex costs at least one by the preceding argument. These costs
use disjoint incidences. A run of length $\ell\ge2$ contains
$\ell-1$ $Z$--$Z$ angles, proving the assertion. Pentagonal facial
boundaries are simple by Lemma~\ref{lem:short-facial-walks}; repeated
occurrences on longer faces do not affect this incidence count.
\end{proof}

Let $T'$ be the heavy vertices other than $w$, and let $a$ count its
degree-four vertices. A clean face at such a vertex $q$ is a pentagon
whose sole heavy vertex is $q$ and which avoids $w$.

\begin{lemma}\label{lem:degree-two-clean-capacity}
Every degree-four vertex in $T'$ has at most one clean face, with
at most one exception. An exceptional vertex has exactly two clean
faces, can occur only if $s=2$, and requires both $x_1,x_2$ to have
degree three. If $\kappa$ counts exceptional vertices and $C$ counts
all clean incidences at degree-four vertices of $T'$, then
\begin{equation}\label{eq:degree-two-clean-slack}
 \kappa=0\ (s\ge3),\qquad 0\le\kappa\le\min\{1,a\}\ (s=2),\qquad
 \delta=a+\kappa-C\ge0.
\end{equation}
\end{lemma}
\begin{proof}
Every noncentral vertex of a clean pentagon has degree three: any
cycle through a member of $S$ contains $w$. A clean pentagon contains
at most one attachment $x_i$. Otherwise two attachments at cycle
distance one or two, together with their degree-two neighbours, give
a regular or one-degree-four shared-hub path of order four or five.
The paths are induced by girth and the endpoint degrees. This also
covers the case that one of the attachments is the centre. Apply
Lemmas~\ref{lem:shared-hub-paths} and~\ref{lem:further-degree-two}.

Two clean faces at $q$ form a hinged pair or a bare/chorded bowtie
and contain at most two attachments. If an attachment $x_i$ is absent,
then $u_i$ is separated from this pair, and
Lemma~\ref{lem:hinged-restorations} or
Lemma~\ref{lem:bowtie-degree-two} applies; its bottom contains $w$.
Thus $s\ge3$ excludes two clean faces. If $s=2$, both attachments
must lie on their union. They must be light: a heavy attachment can
only be the centre, in which case each clean face already contains
it and cannot contain the other attachment.

A light $x_i$ has just one facial angle that can avoid $w$, since its
other two angles use $x_i u_i$ and continue through $w$. Its clean
face is therefore fixed. The two faces containing $x_1,x_2$ fix their
unique heavy centre, so at most one centre has two clean faces. Three
clean faces there would include a pair missing an attachment,
contradicting the same restoration. Summing the capacities proves
the stated slack identity and inequalities.
\end{proof}

For a pentagonal face $F$, let $a(F)$ count its heavy vertices other
than $w$, and let $a_4(F)$ count those of degree four. Let $N_4,N_+$
count the nonpentagonal incidences of vertices in $T'$ of degree four
and at least five, respectively, with multiplicity on facial walks.
Put
\begin{align}
 e&=\sum_{v\in T',\ d(v)\ge6}(d(v)-5),\qquad
 D'=2h'-a+e,\label{eq:degree-two-degree-split}\\
 W&=N_4+2N_+
   +\sum_{|F|=5,\ a(F)\ge2}\bigl(2(a(F)-1)-a_4(F)\bigr).
 \label{eq:degree-two-weighted-slack}
\end{align}
Every summand in $W$ is nonnegative, since
$a_4(F)\le a(F)\le2(a(F)-1)$ when $a(F)\ge2$.

\begin{lemma}\label{lem:degree-two-master}
For $s\ge2$, the quantities above satisfy
\begin{align}
 a+b+\delta+W+2Q+2e&\le2h'-12+\kappa,
 \label{eq:degree-two-master}\\
 D'&\ge12+b-\kappa+2Q+3e+\delta+W,
 \label{eq:degree-two-excess-master}\\
 p&\ge30+2t+2b-2\kappa+5Q+6e+2\delta+2W.
 \label{eq:degree-two-face-master}
\end{align}
\end{lemma}
\begin{proof}
Let $A_4$ count pentagons containing $w$ and exactly one other heavy
vertex, of degree four. Such a pentagon cannot use an $S$--$S$ angle:
its complementary path would be the one-degree-four order-four
reduction. It uses a $Z$-edge incidence not used by a $p_0$-face, so
$A_4\le b$. Partitioning the $4a$ degree-four angles gives
\[
 4a=C+A_4+N_4+\sum_{|F|=5,\ a(F)\ge2}a_4(F).
\]
Combine this identity with \eqref{eq:degree-two-deficit-parts} and
\eqref{eq:degree-two-weighted-slack} to obtain
\[
 3a-\kappa+\delta+W=2R-2b+A_4\le2R-b.
\]
Substitute $R=3h'-D'-6-Q$ and $D'=2h'-a+e$ to get the first two
inequalities. Substitution in $p=6+2t+2D'+Q$ gives the third.
All identities count repeated nonpentagonal incidences with
multiplicity.
\end{proof}

\begin{lemma}\label{lem:degree-two-eight-heavy}
A minimum counterexample containing a degree-two vertex has at least
eight heavy vertices, $D\ge14$, and at least thirty-four distinct
$5$-cycles. If its degree-two vertex is unique, then $D\ge16$ and
it has at least forty-two distinct $5$-cycles.
\end{lemma}
\begin{proof}
The unique case is Lemma~\ref{lem:unique-degree-two-heavy-budget}.
Suppose $s\ge2$. Since $a\ge\kappa$ and $b\ge1$,
\eqref{eq:degree-two-master} forces $h'\ge7$, giving eight heavy
vertices including $w$. For $s\ge3$, use $t\ge1$ and $\kappa=0$ in
\eqref{eq:degree-two-face-master}; for $s=2$, the heavy hub gives
$t\ge2$ and $\kappa\le1$. In both cases $p\ge34$.

Equation~\eqref{eq:degree-two-excess-master} gives $D'\ge13$
unless $s=2$ and $\kappa=1$. Since $D=D'+s+t-3$ and $w$ is heavy,
this already gives $D\ge14$ outside that exceptional case. Suppose
then that $s=2$, $\kappa=1$ and $D\le13$. The same equation and
$t\ge2$ force
\[
 t=2,\quad D'=12,\quad b=1,\qquad Q=e=\delta=W=0.
\]
Both attachments have degree three by
Lemma~\ref{lem:degree-two-clean-capacity}. Since $Q=0$, every face
has a simple boundary of length five or six. An $S$--$S$ pentagon
therefore gives the regular shared-hub path of order four. On an
$S$--$S$ hexagon, $W=0$ excludes any heavy vertex other than $w$,
so the complementary path is the regular shared-hub path of order
five. Both contradict Lemma~\ref{lem:shared-hub-paths}. There is
consequently no $S$--$S$ angle. The cyclic word at $w$ must be
$S,Z,S,Z$, and its two isolated $Z$-runs give $b\ge2$ by
\eqref{eq:degree-two-z-run-cost}, a contradiction. Thus $D\ge14$.
Finally the pentagonal faces give distinct $5$-cycles, as noted
after \eqref{eq:degree-two-euler}.
\end{proof}

\subsection{Two separated five-faces}
\begin{lemma}\label{lem:subcubic-face-count}
Let $X$ be a connected plane graph that contains a cycle, has girth at
least five, and has maximum degree at most three.  If $p$ and $q$ are
the numbers of vertices of degrees one and two, respectively, then
\[
 f_5(X)=12-4p-2q+
 \sum_{\substack{F\text{ a face of }X\\ |F|\ge7}}(|F|-6).
 \tag{FC}
\]
Here $f_5(X)$ counts facial walks of length five.  Every facial walk
has length at least five, and a facial walk of length five is a simple
5-cycle.
\end{lemma}
\begin{proof}
The boundary walk of a face is a closed walk.  Since $X$ is connected
and contains a cycle, each such walk contains a simple cycle; a bridge,
when it occurs on a boundary walk, is traversed twice on the way to and
from the cyclic block.  The girth assumption therefore gives length at
least five.  A non-simple closed walk contains a simple cycle and at
least two further edge traversals, so in this graph it has length at
least seven.  Thus a length-five facial walk is a simple 5-cycle.

Put
\[
 L_X=\sum_{|F|\ge7}(|F|-6).
\]
The face-length sum is $2e=6f-f_5+L_X$, while the degree sum is
$2e=3n-2p-q$.  Combining these equalities with Euler's identity
$n-e+f=2$ gives
\[
 f_5=6f+L_X-2e
     =12-6n+4e+L_X
     =12-4p-2q+L_X,
\]
as claimed.
\end{proof}

\begin{lemma}\label{lem:two-separated-five-faces}
Every connected plane graph of girth at least five with
$2\le d(v)\le3$ and at most one degree-two vertex contains
two vertex-disjoint 5-cycles with no edge between them.
\end{lemma}
\begin{proof}
We use Lemma~\ref{lem:subcubic-face-count} throughout.  Set
\[
 q:=|\{v\in V(G):d_G(v)=2\}|\in\{0,1\},\qquad
 p:=|\{v\in V(G):d_G(v)=1\}|=0.
\]
If $q=1$, let $z$ denote the unique degree-two vertex; when $q=0$
there is no vertex $z$ in use.  Thus $G$ has at least $12-2q$
facial 5-cycles.
If there is a bridge $ab$, take the side $C$ not containing the
unique degree-two vertex, if it exists, and relabel the endpoints so
that $a\in C$.  (If there is no degree-two vertex, choose either
side.)  In the graph induced by this side, deleting $ab$ leaves $a$
of degree two and every other vertex of degree three.  It is cyclic:
otherwise, being connected, it would be a tree, which is impossible
when every vertex has degree at least two.  Lemma~\ref{lem:subcubic-face-count}
gives at least
ten facial 5-cycles.  Since $a$ has two face angles, at most two of
these cycles contain $a$; choose one, $S$, that avoids $a$.

Let $D$ be the other side.  If the original degree-two vertex is in
$D\setminus\{b\}$, then $D$ has two degree-two vertices; if it is
$b$, then $b$ is its only degree-one vertex; and if there is no such
vertex, $D$ has one degree-two vertex.  Thus $4p_D+2q_D\le4$ in all
cases.  The side $D$ is cyclic, since a tree would have at least two
degree-one vertices.  The face-count lemma therefore gives at least
eight facial 5-cycles in $D$.  Any one of them is disjoint from $S$;
it has no edge to $S$ because $ab$ is the only edge between the two
sides and $a\notin S$.

Assume now there is no bridge. The subcubic graph is 2-connected:
at a cutvertex every branch would require at least two incident
edges to avoid a bridge, forcing degree at least four.
Take a facial 5-cycle $S$ and let $B$ be its exterior neighbours.
Girth makes them distinct; $b=|B|$ is five, or four when $q=1$
and $S$ contains the unique degree-two vertex $z$.
Set $H=G-S$, embedded with the deleted region outside.
Every member of $B$ is on that outside face.  A vertex outside $S$
has at most one neighbour on $S$, by the girth assumption, and has
degree at least two in $G$; hence $H$ has no isolated vertex.  Each
member of $B$ has degree two in $H$, except that, if $q=1$ and
$z\in B$, the vertex $z$ becomes a leaf.
There is at most one leaf, so every component of $H$ is cyclic (a
finite tree has at least two leaves).

In one such component, let $r$ count its members of $B$, let $q'$
count the degree-two vertices of $G$ in the component that are not in
$B$, and let $\ell$ be its outside facial length.  Define
\[
 \varepsilon:=
 \begin{cases}
  1,&q=1\text{ and $z\in B$ lies in this component},\\
  0,&\text{otherwise},
 \end{cases}
 \qquad 0\le q'\le1.
\]
When $q=0$, these definitions give $\varepsilon=q'=0$.
The degree-one and degree-two counts in this component are
\[
    p_H=\varepsilon,\qquad
    q_H=r-\varepsilon+q'.
\]
Applying Lemma~\ref{lem:subcubic-face-count} and removing the outside
face gives
\[
 f^{\rm in}_5=\ell+6-2r-2\varepsilon-2q'+L,\qquad
 L=\sum_{\substack{F\text{ internal}\\|F|\ge7}}(|F|-6)\ge0.       \tag{**}
\]
Indeed,
\[
4p_H+2q_H=4\varepsilon+2(r-\varepsilon+q')
              =2r+2\varepsilon+2q'.
\]
The outside
boundary is a closed walk containing a cycle, so $\ell\ge5$.  If
$\ell=5$ it contributes one facial 5-cycle, while if $\ell\ge6$ it
contributes $\ell-6$ to the face-length excess; in either case the
subtraction gives the displayed formula.
A member of $B$ belongs to at most one internal 5-face: its degree
is at most two, with an outside face angle, and a leaf belongs to none.
If $H$ is disconnected, some component has $r\le2$ since
$|B|\le5$.  Since $\varepsilon+q'\le1$ and $\ell\ge5$, (**) gives
\[
 f^{\rm in}_5\ge 5+6-2r-2(\varepsilon+q')\ge9-2r>r.
\]
It therefore has more internal 5-faces than attachment vertices.
One avoids $B$, so it is disjoint from $S$ and has no edge to $S$.

Suppose $H$ is connected.  Let $F_1,\ldots,F_t$ be the distinct
faces on the exterior side of the edges from $S$ to $B$.  Deleting
the five vertices of $S$ and the $5+b$ incident edges gives
$n(H)=n(G)-5$ and $e(H)=e(G)-5-b$; since $H$ is connected, Euler's
formula gives $f(H)=f(G)-b$.  In the deletion, the 5-face bounded by
$S$ and the $t$ exterior faces merge into the new outside face, so
$t=b$.  Thus the exterior faces are indeed distinct; call them
$F_1,\ldots,F_b$.  Their merger has outside length
\[
 \ell=\sum_{i=1}^b|F_i|-(2b+5),                       \tag{***}
\]
because every edge of $S$ occurs once among the $F_i$, while every
edge from $S$ to $B$ occurs twice (once for each incident face).

If $q=0$, any $S$ has $b=5$, $\ell\ge10$, and
$\varepsilon=q'=0$.  Formula (**) gives
\[
 f^{\rm in}_5\ge10+6-2\cdot5=6>|B|.
\]
If $q=1$ and $z$ lies on a 5-face, choose that face as $S$. Now $b=4$,
$\ell\ge7$, and $\varepsilon=q'=0$ (the unique degree-two vertex is
on $S$, not in $B$).  Hence
\[
 f^{\rm in}_5\ge7+6-2\cdot4=5>|B|.
\]

It remains that $q=1$ and $z$ lies on no 5-face. There is a 5-face adjacent
to two distinct non-5-faces. To see this, consider the plane dual.
It is loopless because $G$ is bridgeless.  Since $G$ is 2-connected,
its facial boundaries are cycles.  Thus the faces of the dual are
triangles, one for each degree-three vertex of $G$, except for the
unique digon corresponding to $z$.  The digon's endpoints are
distinct (there is no cutvertex) and represent the two faces incident
with $z$; both are
non-5-faces because $z$ lies on no 5-face.  Let $W$ be the set of
dual vertices representing non-5-faces.  Suppose, for a contradiction,
that every dual vertex outside $W$ has at most one distinct neighbour
in $W$.  Then in every triangular dual face whose two consecutive
vertices lie in $W$, the third vertex also lies in $W$; otherwise it
would have two distinct neighbours in $W$.

Start with the digon edge.  Its endpoints lie in $W$.  Walk around
either endpoint in the fixed plane embedding.  Each intervening face
is triangular, except possibly the unique digon, so the preceding
paragraph shows inductively that every successive neighbour is in
$W$; the digon only repeats the same neighbour.  The same argument
applies at every newly reached vertex of $W$.  Hence $W$ is closed
under adjacency.  The dual is connected, so this would give
$V(G^*)=W$, contradicting the at least ten 5-faces of $G$, whose dual
vertices are not in $W$.  Therefore some 5-face is adjacent to two
distinct non-5-faces, as asserted.
Choose the resulting 5-face $S$. If its deletion disconnects the graph,
use the case already proved; otherwise $b=5$ and two distinct
neighbours in (***) have length at least six.  Thus $\ell\ge12$.
If $q=1$ and $z\in B$, then $\varepsilon=1,q'=0$; if
$q=1$ and $z\notin B$, then $\varepsilon=0,q'\le1$.  In both cases
$\varepsilon+q'\le1$, so (**)
gives
\[
 f^{\rm in}_5\ge12+6-2\cdot5-2=6>|B|.
\]
In every case one avoids $B$, making it disjoint and nonadjacent to $S$.
\end{proof}

\paragraph{Completing the subcubic induction.}
\begin{proof}[Proof of Theorem~\ref{thm:subcubic}]
Induct on the number of vertices, allowing disconnected graphs.
The one-vertex case is $\R_4(G)=K_4$. A vertex of degree at most one
is handled by Lemma~\ref{lem:single-lift}.
If there are two degree-two vertices, adjacent ones use Lemma~\ref{lem:adjacent-two-four}. Ones at distance two use Lemma~\ref{lem:path232}, unless their
middle vertex has degree two, when an adjacent pair already applies.
Ones at distance at least three use Lemma~\ref{lem:separated-two-g5}, whose spare-colour witness uses
Gr\"otzsch's theorem~\cite{Grotzsch59}; the new pentagon lift does not.
Pairs in different components have disjoint neighbourhoods and
are covered by the same separated-pair argument.

The remaining graph has minimum degree at least two and at most
one degree-two vertex. Apply Lemma~\ref{lem:two-separated-five-faces}
inside any component. Delete its two separated pentagons.
The remaining graph is nonempty and belongs to the same class, so
induction gives its four-colour Hamilton cycle.
Lemma~\ref{lem:two-pentagons-restore} restores both pentagons.
This proves the assertion at four colours.
For every $k\ge5$ use the triangle-free planar theorem proved earlier,
with larger colours separately supplied there by degeneracy.
The frozen $012012$ colouring of $C_6$ gives the Hamiltonian
threshold lower bound four.
\end{proof}


\subsection{Graphs with subcubic blocks}
\begin{theorem}\label{thm:block-subcubic}
Let $G$ be a nonempty finite simple planar graph of girth at least
five. If every block $B$ satisfies $\Delta(B)\le3$, then $\Rk(G)$ is
Hamiltonian for every $k\ge4$. Cutvertices may have arbitrarily large
degree in $G$. The uniform Hamiltonian threshold of this class is exactly four.
Its Hamilton-connected threshold lies between four and five.
\end{theorem}
\begin{proof}
Treat bridges as two-vertex blocks and isolated vertices as singleton
blocks. The class is closed under vertex deletion: a nontrivial
two-connected block of a subgraph is a two-connected subgraph of
$G$, hence is contained in one original block, with no larger
internal degrees. Planarity and girth are also preserved.

Products of Hamiltonian graphs are Hamiltonian: copy one factor
cycle along a spanning path of another factor cycle, and merge
successive copies using two alternating, distinct internal edges.
This also handles disconnected original graphs. Choose a smallest
four-colour counterexample, which may therefore be assumed connected.
The singleton base is $K_4$, and the single-vertex lift excludes
vertices of degree at most one. If there is only one block, the
previous subcubic girth-five theorem applies. Thus assume minimum
degree at least two and at least one cutvertex.

Lemmas~\ref{lem:adjacent-two-four}, \ref{lem:path232}, and
\ref{lem:separated-two-g5} exclude adjacent degree-two vertices, a
$(2,3,2)$ path, and two degree-two vertices at distance at least three.
Consequently any two degree-two vertices not already giving a
reduction must have a common neighbour of degree at least four in $G$.

The block-cut tree has at least two endblocks. No endblock is a
bridge, since its noncut endpoint would be a leaf. Let $B$ be an
endblock and $x$ its unique cutvertex of $G$. Every vertex of
$B-x$ has its full $G$-degree in $B$, hence degree two or three;
also $d_B(x)\in\{2,3\}$. Let $q$ count the degree-two vertices of $B-x$ and put
$\varepsilon=[d_B(x)=2]$. If $q\ge2$, any pair of these vertices
has a common neighbour of $G$-degree at least four by the excluded
reductions. Their neighbours all lie in $B$, so only $x$ can be
this common neighbour. Every one of the $q$ vertices is therefore
adjacent to $x$, giving $q\le d_B(x)=3-\varepsilon$. Together with
the immediate $q\le1$ case, this proves $q+\varepsilon\le3$.

Euler's identity in the two-connected block $B$ now gives
\[
 f_5(B)=12-2(q+\varepsilon)+\sum_{|F|\ge7}(|F|-6)\ge6.
\]
Each five-face is a simple pentagon, while at most $d_B(x)\le3$
faces meet $x$, so choose one avoiding $x$. If $q\le1$, it contains
at most one degree-two vertex of $G$. If $q\ge2$, a cycle through
any of the $q$ vertices must use its edge to $x$, so the chosen
pentagon contains none. All its other vertices have $G$-degree three.

Choose these pentagons in two distinct endblocks. Their vertices
are noncutvertices of $G$ and belong to their respective endblocks
only, so the pentagons are disjoint and have no edge between them.
The established protected-edge lift for two separated pentagons
applies, including its allowance of one degree-two vertex per
pentagon and simultaneous exterior list changes. The smaller graph
remains in the class and has a bottom Hamilton cycle by induction.
Restoration contradicts minimality.

For every $k\ge5$ separately use the established triangle-free
planar theorem. Finally $C_6$ belongs to the class and its periodic
$012012$ colouring is frozen in $\R_3(C_6)$, giving the lower bound
four.
\end{proof}
In particular, a remaining minimum counterexample for general
planar girth-five four-colour Hamiltonicity must contain a block
with a vertex of degree at least four \emph{inside that block}.
High total degree accumulated only at cutvertices is insufficient.
This does not assert that every minimum counterexample is two-connected.


\subsection{Small total degree excess}
Recall that $D(G)$ is the total degree excess above three.
It cannot increase under vertex deletion. All degree conditions
in the following structural arguments refer to the original plane graph.

\begin{lemma}\label{lem:hinged-neighbourhood}
Let $G$ be a connected plane graph of minimum degree at least three and
girth at least five. Suppose a hinged pair $X$ consists of two facial
pentagons. Put $B=N(X)\setminus X$ and
$D_B=\sum_{v\in B}(d(v)-3)$. At most $15+D_B$ pentagonal faces meet
$X\cup B$.
\end{lemma}
\begin{proof}
Use the local labels $(0,2,3,4,1,0)$ and $(0,1,7,6,5,0)$ for the
pentagons, with sole heavy vertex $0$. Their union is a disk. In the
bare model its exterior boundary has seven outside incidences, at
seven different local vertices. The sectors between them account for
at most seven further faces meeting $X$, hence at most nine pentagonal
faces in total. Bridges may merge sectors and only reduce this count.
In the chorded model the chord $36$ divides the exterior into two
regions and leaves five outside incidences. A region with contacts
has at most as many sector faces as contacts; a contact-free region
has at most one. At most one region is contact-free, giving at most
eight pentagonal faces meeting $X$.

Let $c\in\{0,1\}$ indicate the chord, put $r=7-2c$, and let
$j=r-|B|$. At $v\in B$ write $r_v=|N(v)\cap X|$. The permitted
outside-role identifications give $r_v\le2<d(v)$. At least $r_v+1$
angles in its circular rotation use an edge into $X$. Thus at most
$d(v)-r_v-1$ angles can belong to a pentagon avoiding $X$. Summing
over $B$ bounds the additional pentagonal faces by
$2|B|-r+D_B$. Consequently
\[
 p_{\rm near}\le16-3c-2j+D_B.
\]
This proves the desired bound unless $c=j=0$.

In that remaining case all seven outside neighbours are distinct.
The port-free endpoint $1$ gives the exterior sector $4,1,7$. If
its face is not pentagonal, the bound of nine loses one. Otherwise
that face is $(4,1,7,b_7,b_4,4)$, where $b_4,b_7\in B$ are distinct,
so $b_4b_7$ is an edge on a cycle. Its other incident face is different.
At both endpoints its angle on that side uses no edge into $X$:
the endpoint has degree at least three and just one neighbour in $X$.
These are two angles included in the preceding candidate sum. If the
other face is a pentagon avoiding $X$, they count the same face twice.
If it meets $X$, it was already counted there; if it is nonpentagonal,
neither angle contributes a new pentagon. In every case at least one
unit is saved, proving $p_{\rm near}\le15+D_B$.
\end{proof}

\begin{lemma}\label{lem:unique-ordinary-partner}
Let $G$ be a connected plane graph of minimum degree at least three,
girth at least five, and $D(G)\le11$. If it has exactly one ordinary
facial pentagon, that pentagon is separated from a hinged pair or bowtie.
\end{lemma}
\begin{proof}
Let $a,b$ count vertices of degrees four and at least five, respectively,
and put $e=\sum_{d(v)\ge6}(d(v)-5)$. Thus $h=a+b$ and
$D=a+2b+e$. If $S$ is the unique ordinary facial pentagon, the face
identity and heavy-incidence count give
\[
 p=12+2D+Q,\qquad
 \Delta=3h+D-(p-1)=2a+b-e-11-Q\ge0.
\]
The deficit $\Delta$ counts heavy incidences on nonpentagonal faces
and excess heavy incidences on pentagonal faces.

Suppose there is no stated partner. Delete $S$ and let $B$ be its
five distinct exterior neighbours. Every residual component $C$ has
minimum degree at least two. The connected set consisting of $S$ and
the other components lies in one face of $C$; choose it as outside.
All attachments lie on that face. Let $r=|B\cap V(C)|$, let $\ell$
be its outside facial length, and let $L$ be its internal face excess
above six. Use $a_C,b_C,e_C,D_C$ for its original degree counts and
excess, and let $\beta_+$ count its attachments of original degree at
least five. Its degree and face sums give the internal pentagon count
\[
 I_C=\ell+6-2r+2D_C+L.
\]
Its internal faces are unchanged faces of $G$. Let $\Delta_C$ be
their contribution to the heavy-incidence deficit, so
$\sum_C\Delta_C\le\Delta$.

An attachment of original degree $d$ has at most $d-2$ internal
angles. Thus at most $r+D_{B\cap V(C)}$ internal pentagons meet
attachments. None of the other internal pentagons is ordinary. A
nonattached degree-four vertex is sole heavy on at most one of them:
two facial pentagons at that vertex form a hinged pair or bowtie
separated from $S$. At a higher-degree vertex use capacity $d$.
Multi-heavy pentagons consume at least one unit of $\Delta_C$ each.
The total sole-heavy capacity is $a_C+5b_C+e_C$. Removing the attached
capacities cancels $D_{B\cap V(C)}$ at degree four and leaves $-3$
per higher-degree attachment. Hence
\[
 I_C\le r+a_C+5b_C+e_C-3\beta_++\Delta_C,
\]
or
\begin{equation}\label{eq:unique-ordinary-component}
 \ell+6-3r+a_C-b_C+e_C+L+3\beta_+\le\Delta_C.
\end{equation}

Every residual component is cyclic, so $\ell\ge5$. If there are
$c\ge2$ components, their attachment counts sum to five. Summing
\eqref{eq:unique-ordinary-component} and using the displayed value
of $\Delta$ gives
\[
 D\ge11c-4+Q+3e\ge18,
\]
a contradiction. Thus $G-V(S)$ is connected. Deleting five vertices
and ten edges reduces the face count by five; the incident faces merge
into one, so the five exterior faces beside $S$ are distinct. If
$\epsilon$ is their total excess above five, their merger has outside
length $10+\epsilon$. Substitution in
\eqref{eq:unique-ordinary-component}, now with $r=5$, gives
\[
 D\ge12+Q+3e+\epsilon+L+3\beta_+\ge12,
\]
again contradicting $D\le11$. This argument permits bridges and
cutvertices throughout.
\end{proof}

\begin{lemma}\label{lem:small-excess-structure}
Let $G$ be a connected plane graph of minimum degree at least three,
girth at least five, and $D(G)\le9$. It has two ordinary facial
pentagons, an ordinary facial pentagon separated from a hinged pair
or bowtie, or a hinged pair separated from a hinged pair or bowtie.
\end{lemma}
\begin{proof}
Two ordinary faces give the first alternative; a unique one is handled
by Lemma~\ref{lem:unique-ordinary-partner}. Suppose there are none.
Retain $a,b,e,h,D,Q$ from the preceding proof. Then
\[
 p=12+2D+Q,\qquad \Delta=3h-D-12-Q\ge0.
\]
The nonpentagonal and multi-heavy pentagonal incidences at degree-four
vertices total at most $2\Delta$: a pentagon with $k\ge2$ heavy
vertices has at most $k\le2(k-1)$ such incidences. The number $C_4$
of sole-heavy pentagonal incidences at degree-four vertices therefore
satisfies $C_4\ge4a-2\Delta$. Since
\[
 a-\Delta=12-h+e+Q>0
\]
by $h\le D\le9$, we have $C_4>2a$. Some degree-four vertex has at
least three sole-heavy pentagonal faces, including two adjacent in
its rotation. Their union is a hinged pair $X$.

Put $B=N(X)\setminus X$ and call a pentagon far if it avoids $X\cup B$.
Suppose $X$ has no separated hinged or bowtie partner. A degree-four
vertex outside $X\cup B$ is sole heavy on at most one far pentagon,
since two would form such a partner. A higher-degree vertex has
capacity $d$. The total sole-heavy capacity is
\[
 C=a+\sum_{d(v)\ge5}d(v)=3h+D-3a.
\]
Remove the capacity one of the heavy vertex of $X$, and the capacities
$C_B$ at heavy vertices of $B$. At most $\Delta$ far pentagons are
multi-heavy. Lemma~\ref{lem:hinged-neighbourhood} consequently gives
\[
 p\le15+D_B+C-1-C_B+\Delta.
\]
If $\beta_+$ counts vertices of $B$ of degree at least five, then
$D_B-C_B=-3\beta_+$. Subtracting the asserted upper bound from the
actual value of $p$ yields
\[
 p-\bigl(15+D_B+C-1-C_B+\Delta\bigr)
   =10-D+3e+2Q+3\beta_+>0,
\]
a contradiction. Thus the last alternative holds.
\end{proof}

\begin{theorem}\label{thm:excess-nine}
For every nonempty planar graph $G$ of girth at least five with
$D(G)\le9$, the graph $\Rk(G)$ is Hamiltonian for every $k\ge4$.
This class has uniform Hamiltonian threshold four.
\end{theorem}
\begin{proof}
The class is closed under vertex deletion. A minimum counterexample
at four colours is connected and has minimum degree at least two by
Lemma~\ref{lem:degree-two-counterexample}, including the singleton
base case. Lemma~\ref{lem:degree-two-eight-heavy} excludes degree two,
since it would require $D\ge14$.
Lemma~\ref{lem:small-excess-structure} supplies a deletion restored by
Lemma~\ref{lem:any-two-ordinary}, Lemma~\ref{lem:hinged-restorations},
or Theorem~\ref{thm:pentagon-bowtie-restoration}. Each deletes at most
seventeen vertices, whereas minimum degree three and planar girth five
give $|V(G)|\ge20$. The bottom graph is nonempty and remains in the
class. Induction gives its Hamilton cycle and hence a Hamilton path,
so restoration contradicts minimality. For $k\ge5$, apply
Theorem~\ref{thm:trianglefree-hc}. The graph $C_6$ belongs to the class
and its three-colouring $012012$ is frozen, proving sharpness.
\end{proof}

\begin{theorem}\label{thm:seven-heavy}
For every nonempty planar graph $G$ of girth at least five with at
most seven heavy vertices, $\Rk(G)$ is Hamiltonian for every $k\ge4$.
This class has uniform Hamiltonian threshold four.
\end{theorem}
\begin{proof}
The heavy-count class is hereditary. A minimum four-colour
counterexample is connected with minimum degree at least two;
Lemma~\ref{lem:degree-two-eight-heavy} excludes degree two.
If $D\le9$, Theorem~\ref{thm:excess-nine} applies. Otherwise let
$o$ count ordinary facial pentagons. The incidence bound is
\[
 o\ge12-3h+D+Q\ge1.
\]
It is at least two unless $h=7$, $D=10$, and $Q=0$. Two ordinary
faces are restored by Lemma~\ref{lem:any-two-ordinary}. If there is
just one, Lemma~\ref{lem:unique-ordinary-partner} supplies a separated
hinged pair or bowtie, since $D=10\le11$. Use
Lemma~\ref{lem:hinged-restorations} or
Theorem~\ref{thm:pentagon-bowtie-restoration}. These deletions have
at most fourteen vertices, so the bottom is nonempty by $|V(G)|\ge20$.
It remains in the heavy-count class and has a Hamilton cycle by
induction, contradicting restoration. Larger palettes follow from
Theorem~\ref{thm:trianglefree-hc}, and $C_6$ gives sharpness.
\end{proof}

\begin{theorem}\label{thm:thirty-one-pentagons}
For every nonempty planar graph $G$ of girth at least five with at
most thirty-one distinct $5$-cycles, $\Rk(G)$ is Hamiltonian for
every $k\ge4$. This class has uniform Hamiltonian threshold four.
\end{theorem}
\begin{proof}
The class is hereditary. A minimum four-colour counterexample is
connected with minimum degree at least two.
Lemma~\ref{lem:degree-two-eight-heavy} excludes degree two, which
would require at least thirty-four distinct $5$-cycles. In minimum
degree at least three, distinct pentagonal faces give distinct
$5$-cycles. Hence
\[
 12+2D+Q=p\le31,
\]
so $D\le9$, contrary to Theorem~\ref{thm:excess-nine}. Larger palettes
follow from Theorem~\ref{thm:trianglefree-hc}; the frozen colouring of
$C_6$ gives sharpness.
\end{proof}

\begin{corollary}\label{cor:two-degenerate-thirty-three}
Every nonempty two-degenerate planar graph $G$ of girth at least five
has Hamiltonian $\Rk(G)$ for every $k\ge4$ if it has at most
thirty-three distinct $5$-cycles or satisfies $D(G)\le13$.
Each of these two classes has uniform Hamiltonian threshold four.
\end{corollary}
\begin{proof}
Both classes are hereditary. A minimum four-colour counterexample
is connected with minimum degree at least two, and two-degeneracy
forces a degree-two vertex in the graph itself.
Lemma~\ref{lem:degree-two-eight-heavy} then requires at least
thirty-four distinct $5$-cycles and $D\ge14$, contradicting the
respective class bound. For $k\ge5$, use
Theorem~\ref{thm:degenerate-hc} with degeneracy two. The graph $C_6$
belongs to both classes and gives sharpness.
\end{proof}

\section{The girth-five frontier}\label{sec:sharpness}

The threshold intervals in Table~\ref{tab:thresholds} leave two
questions: does every $G\in\mathcal P_5$ have a Hamiltonian
$\R_4(G)$, and is every such recolouring graph Hamilton-connected?
The latter question also remains open for the four-colour subclasses
of Theorem~\ref{thm:main-four}. Known connectivity~\cite{BBFHMP23}
does not settle either question. The new subclass theorems require
any minimum Hamiltonian counterexample to have at least thirty-two
distinct $5$-cycles, at least eight heavy vertices, and $D(G)\ge10$.
Theorem~\ref{thm:block-subcubic} also requires a block with internal
maximum degree at least four. If the counterexample contains a
degree-two vertex, Lemma~\ref{lem:degree-two-eight-heavy} gives at
least thirty-four distinct $5$-cycles and $D(G)\ge14$; if that vertex
is unique, these bounds become forty-two and sixteen.

\begin{remark}\label{rem:thirty-three-frontier}
Suppose a minimum counterexample in the full class $\mathcal P_5$
has at most thirty-three distinct $5$-cycles. The degree-two bound
gives minimum degree at least three. Euler's identity and
Theorem~\ref{thm:excess-nine} then force $D=10$ and $Q\le1$.
Two ordinary facial pentagons are reducible, and a unique one is
excluded by Lemma~\ref{lem:unique-ordinary-partner}, since $D\le11$.
Hence there is none.

The incidence argument in Lemma~\ref{lem:small-excess-structure}
still gives a hinged pair, since $h\le D=10<12$. A separated partner
would restore, so the same neighbourhood calculation now requires
\[
 10-D+3e+2Q+3\beta_+\le0.
\]
Consequently $e=Q=0$. The face count is $p=32$, every face is a
pentagon or hexagon, and at least eight vertices are heavy. The
heavy-degree multiset is therefore six fours and two fives, eight
fours and one five, or ten fours; all other degrees are three.
These are necessary conditions, with no counterexample asserted.
Here $32$ counts pentagonal faces, while the hypothesis counts all
distinct $5$-cycles.
\end{remark}

\section{Finite certificates and reproducibility}
\label{sec:verification}

The supplementary source \texttt{article\_certificates.tex} supplies the
adjacent-pair, $(2,3,2)$ path and pentagon data. The literal cycles for
paths with a common outside neighbour are supplied in
\texttt{path4-family.json} and \texttt{path5-family.json} under
\texttt{supplement/data/}. The seven further degree-two families are in
\texttt{supplement/data/degree-two/}. Three pentagonal families are in
\texttt{supplement/data/pentagonal-reservoirs/}. Both sets have separate
direct checks of their selected-cycle interfaces in physical colour labels.
The extended cycles and bowtie seeds and clique trees are in
\texttt{supplement/data/extended-cycles/} and
\texttt{supplement/data/bowties/}. Hinged-pair product seeds and clique
trees are in \texttt{supplement/data/hinged-pairs/}.
The fourteen exact
readers reconstruct the local graphs and their stated interfaces.
Missing blocks, duplicate records, illegal transitions,
incomplete state coverage and failed protection conditions cause the
check to fail. No solver output or earlier computation report is an
input to these checks.

\begin{table}[htbp]
\centering\small
\caption{Finite checks supporting the local restoration lemmas.}
\label{tab:verification}
\begin{tabularx}{\linewidth}{@{}lX@{}}
\toprule
Supplement & Checked objects and coverage\\
\midrule
Appendix A & Seven edge-covering cycle orbits; 53 endpoint words with
exact subset-DP cross-check; 225 labelled list pairs and 3,225 interfaces.\\
Appendix B & Three normalized path cycles; 60 good boundaries and all
504 directed adjacent good-boundary changes; safe edges for every
allowed cycle pair and support at bad neighbours.\\
Appendices C--D & 23 complete pentagon components; 2,304 labelled
boundaries; ordered path blocks; both hard two- and three-row seeds;
pool membership and protection in every row; all 313 interfaces,
partitioned as $47+236+30$; all $23^2$ product cycles with the
alternating reserve splice replay; and the relative-normalization rule
on all labelled boundaries and admissible maps (75,120 overlapping
cases).\\
Appendix E & Five order-four and fifteen order-five path cycles;
all 64 and 256 labelled boundaries and all 576 and 3,072 directed
physical-boundary interfaces; at least three and five common chosen-cycle
edges, respectively.\\
Appendix F & Seven degree-two cycle families with 421 templates and
96,036 template states; all 8,704 labelled boundaries and 158,976 directed
interfaces, including every allowed paired-port move; a separate physical
replay of all interface inequalities.\\
Appendix G & Three pentagonal families with 781 templates, 786 component
cycles and 101,121 template states; all 17,664 labelled boundaries and
486,912 directed interfaces. Every overlapping component pair shares at
least four selected-cycle edges in the reservoir and at least two in the
other models; every component has an overlapping partner on both sides.
A separate physical replay checks all interface inequalities.\\
Appendix H & Five connected-cycle families with 771 templates and
2,491,162 states; all 75,008 labelled boundaries and 2,577,024 directed
interfaces, with complete alias matchings and independent physical
histograms.\\
Appendix I & 558 bowtie fibres with 845,440 states; 1,116 $K_2/K_3$
seed cycles; all 2,251,776 directed factor interfaces before and after
reserves; complete clique trees of physical capacity two; and the
ordinary-pentagon seed extract bound to the marked data.\\
Appendix J & 410 hinged-pair fibres with 267,825 states; 1,230
$K_2/K_3/P_3$ seed cycles and 446,976 directed full/core interfaces;
complete clique trees of physical capacity two; independent seed
coverage and literal physical interface histograms.\\
\bottomrule
\end{tabularx}
\end{table}

Run \texttt{python3 supplement/verify\_all.py} from the article directory.
For the standalone delivery directory, add
\texttt{--check-manifest} to verify every listed file byte for byte.
The scripts require only the Python standard library and use exact
integer and finite-set operations. The runner also checks rejection
of corrupted cycles, missing certificate blocks, invalid interface
edges and altered selector data. Its JSON reports record input hashes
and coverage; the distribution includes commands
for rebuilding both PDFs and a file manifest.

The finite verification establishes the local certificate claims, the
finite $23^2$ product/reserve replay and the relative-normalization
replay in Table~\ref{tab:verification}.
The path lifting, relative normalization and component-graph
connectivity along an arbitrary bottom Hamilton path, spanning-tree splicing and structural inductions are
analytic arguments in the article. Thus the finite checks have no bound
on the order of $G$ as an input premise, and the general girth-five
questions stated above remain open.

\Needspace{20\baselineskip}
\begin{thebibliography}{9}
\bibitem{BBFHMP23}
V.~Bartier, N.~Bousquet, C.~Feghali, M.~Heinrich, B.~Moore, and T.~Pierron,
\newblock Recolouring planar graphs of girth at least five,
\newblock \emph{SIAM Journal on Discrete Mathematics} 37 (2023), no.~1, 332--350.
\newblock \doi{10.1137/21M1463598}.

\bibitem{BP16}
N.~Bousquet and G.~Perarnau,
\newblock Fast recoloring of sparse graphs,
\newblock \emph{European Journal of Combinatorics} 52A (2016), 1--11.
\newblock \doi{10.1016/j.ejc.2015.08.001}.

\bibitem{ChMa11}
K.~Choo and G.~MacGillivray,
\newblock Gray code numbers for graphs,
\newblock \emph{Ars Mathematica Contemporanea} 4 (2011), no.~1, 125--139.
\newblock \doi{10.26493/1855-3974.196.0df}.

\bibitem{Grotzsch59}
H.~Gr\"otzsch,
\newblock Zur Theorie der diskreten Gebilde, VII: Ein Dreifarbensatz f\"ur dreikreisfreie Netze auf der Kugel,
\newblock \emph{Wiss. Z. Martin-Luther-Univ. Halle-Wittenberg, Math.-Nat. Reihe} 8 (1958/1959), 109--120.

\bibitem{HLTH04}
H.-C.~Hsu, T.-K.~Li, J.~J.~M.~Tan, and L.-H.~Hsu,
\newblock Fault Hamiltonicity and fault Hamiltonian connectivity of the arrangement graphs,
\newblock \emph{IEEE Transactions on Computers} 53 (2004), no.~1, 39--53.
\newblock \doi{10.1109/TC.2004.1255789}.

\end{thebibliography}

\end{document}
